% appendix.tex
% Supplementary material: empirical validations of assumptions used in
% the main-text derivations, plus material that is too dense for the
% main body but useful for replication.
%
% Language-audit pass, 2026-07-14 (see project_language_audit_2026-07-14
% memory; final section of the pass). Kept every algebra step (no
% collapsed align blocks) but cut narrating prose around them ("we
% expand this in two stages", "distributing", "pulling outside the
% sums"). Killed generic terminology: "artifact" -> "bias" (matches
% discussion.tex's "diversification-return bias"; the one surviving
% "sample-size artifact" is a standard statistical idiom, kept
% consistent with results.tex), "anchor"/"anchoring" -> literal
% ("set from", "matching"), "engine"/"engine-raw" -> "raw" (was
% implementation jargon for the uncorrected ensemble; also renamed the
% $W_{\mathrm{engine}}$ subscript to $W_{\mathrm{raw}}$), "operating
% point" -> "parameter values"/"fitted", "pipeline"/"framework" ->
% "construction"/"procedure", generic "mechanism" -> literal (e.g. "the
% persistence mechanism" -> "the duration model"), "downstream" cut,
% "smoke test" cut (software-testing jargon), "composer" -> "composition
% method" (last file needing this swap; done everywhere now), "sat/sits"
% -> "clustered"/"remained near"/"falls". Table 3 caption's "Clipping is
% a stability mechanism" rewritten to state the observed result
% directly, per the audit's own suggested rewrite.

\section{Empirical cross-term diagnostic}
\label{sec:supp-cov}

The variance correction of Eq.~\eqref{eq:scale} assumes
$\Cov(\tg_i^{\mathrm c},g_m)\approx0$, allowing the variance derivation to omit
the cross-term. Each generator draw is simulated independently of the
market path, so this covariance is zero in population, but on a
finite path its sample value is not. We measured the size of the
sample cross-term for three residual sources: the full-return
JumpHMM generator used by the hybrid and naive composition methods,
the residual-fit JumpHMM generator used by the JumpHMM-on-residuals
method, and the Gaussian residual draw of the Gaussian SIM. For each
of the $423$ non-market tickers, we sampled $R = 100$ training-length
growth-rate paths $\tg_i^{(r)}$ from each source and computed the
sample correlation of each centered draw with the observed training
growth-rate path $g_m$ of SPY, averaged over the $R$ draws:
\begin{equation}
  \mathrm{cov}^{\mathrm{norm}}_i
  = \frac{1}{R} \sum_{r=1}^{R}
    \frac{\Cov((\tg_i^{(r)})^{\mathrm c},\, g_m)}{\sigma_{\gen,i}^{(r)}\, \sigma_m},
\end{equation}
where $\sigma_{\gen,i}^{(r)}$ is the sample standard deviation of
draw $r$. Centering does not change covariance, so these values equal
those obtained from the raw draws.
We plotted the per-ticker values as histograms
(Fig.~\ref{fig:cov_diag}). The universe means were
$1.53\times10^{-4}$, $6.58\times10^{-5}$, and $2.26\times10^{-4}$ for
the full-return, residual-fit, and Gaussian sources, respectively,
and the per-ticker values of all three ranged from $-0.0061$ to
$0.0061$. For two independent series of $2{,}766$ observations, the
sample correlation has standard deviation near $1/\sqrt{2{,}766}
\approx 0.019$, the spread we observed across single draws, so an
average over $100$ draws has standard deviation near $0.002$. The
per-ticker values are therefore of the size expected from sampling
noise, and the omitted cross-term is negligible on average. On a
single path, it is not zero, which is why the composed variance
matches its target only approximately.

\begin{figure}[H]
  \centering
  \includegraphics[width=0.85\textwidth]{figs/supplement/FigS01-Cross-Covariance.pdf}
  \caption{\textbf{Empirical normalized cross-covariance
    $\Cov(\tg_i^{\mathrm c}, g_m) / (\sigma_{\gen,i}\, \sigma_m)$ across
    the $423$ non-market tickers.} Each per-ticker value is the sample
    correlation between a centered generator draw and the observed
    training SPY growth-rate path, averaged over $R = 100$ draws;
    the histograms for the three residual sources are stacked. Universe means were $1.53\times10^{-4}$,
    $6.58\times10^{-5}$, and $2.26\times10^{-4}$ for the full-return
    JumpHMM, residual-fit JumpHMM, and Gaussian generators.}
  \label{fig:cov_diag}
\end{figure}

% =============================================================================
% S2  Bias correction for daily-rebalanced backtests
% Four flat paragraphs:
%   (1) origin of the bias (Booth-Fama identity + iid residuals)
%   (2) empirical demonstration on a 20-asset allocator
%   (3) the location-shift correction (formula, prior, validation check)
%   (4) rejected alternatives
% =============================================================================

\section{Rebalancing under independent residuals}
\label{sec:supp-bias}

To determine what a rebalanced portfolio of composed paths
(Eq.~\eqref{eq:sim-composition}) inherits from the independent-residual
construction, this appendix derives the portfolio's diversification
return. On the variance-preserving branch, the variance correction in
Eq.~\eqref{eq:scale} sets each composed variance to the generator
variance. When the generator variance matches the empirical variance,
the per-asset variance terms of the diversification return derived
below are matched as well. What the construction does not match is
the cross-asset residual covariance, and the derivation isolates where
that omission enters.

The composition is built to reproduce per-asset marginals and
single-factor cross-sectional dependence. Each asset's residual path,
whether a rescaled full-return draw or a draw from a residual model,
keeps its own temporal structure, but the residual paths are generated
independently across assets at each date, so the construction omits
cross-asset residual covariance. Consider a daily-rebalanced portfolio
with constant target weights $w_i \ge 0$, $\sum_i w_i = 1$,
reset to those weights at the start of each trading day. Throughout
this derivation $g_i(t)$ is the annualized log-growth rate of
Eq.~\eqref{eq:growth_rate}, so the per-period log return on asset
$i$ is $g_i(t)\,\Delta t$ with $\Delta t = 1/252$~yr. The wealth
process is:
\begin{equation}
  W_t \;=\; W_{t-1} \sum_{i} w_i \exp\!\big(g_i(t)\,\Delta t\big),
  \label{eq:wealth-recursion}
\end{equation}
so the per-step log return on the portfolio is
$\log(W_t/W_{t-1}) = \log\!\sum_i w_i \exp(g_i(t)\,\Delta t)$. We
suppress the time argument, writing $g_i$ for $g_i(t)$. A
second-order Taylor expansion of the exponential about zero is
$\exp(g_i\,\Delta t) = 1 + g_i\,\Delta t + \tfrac{1}{2}(g_i\,\Delta t)^2
+ O((g\,\Delta t)^3)$. Substituting it into each summand of
$\sum_i w_i \exp(g_i\,\Delta t)$ gives:
\begin{align}
  \sum_{i} w_i \exp\!\big(g_i\,\Delta t\big)
   &= \sum_{i} w_i \!\left[\,1 + g_i\,\Delta t
        + \tfrac{1}{2}(g_i\,\Delta t)^2
        + O\!\big((g\,\Delta t)^3\big)\right]
       && \text{(Taylor)} \notag \\
   &= \sum_{i} w_i
        + \Delta t \sum_{i} w_i\, g_i
        + \tfrac{1}{2}(\Delta t)^2 \sum_{i} w_i\, g_i^2
        + O\!\big((g\,\Delta t)^3\big)
       && \text{(distribute)} \notag \\
   &= 1 + \Delta t \sum_{i} w_i\, g_i
        + \tfrac{1}{2}(\Delta t)^2 \sum_{i} w_i\, g_i^2
        + O\!\big((g\,\Delta t)^3\big).
       && \text{(}\textstyle\sum_i w_i = 1\text{)}
  \label{eq:exp-expand}
\end{align}
A second-order Taylor expansion of $\log(1+x)$ about $x=0$ gives
$\log(1+x) = x - \tfrac{1}{2}x^2 + O(x^3)$. Identifying $1 + x$ with
the right-hand side of Eq.~\eqref{eq:exp-expand}, the part beyond one
is
$x :=\Delta t \sum_i w_i\, g_i + \tfrac{1}{2}(\Delta t)^2 \sum_i w_i\, g_i^2$,
and squaring keeps only the leading $(\Delta t)^2$ term:
$x^2 = (\Delta t)^2 (\sum_i w_i\, g_i)^2 + O((\Delta t)^3)$ because
the cross term has an extra $\Delta t$ factor. Substitution gives:
\begin{align}
  \log\!\frac{W_t}{W_{t-1}}
   &= \log\!\Big(1 + x + O\!\big((g\,\Delta t)^3\big)\Big) \notag \\
   &= x - \tfrac{1}{2}\, x^2 + O\!\big(x^3\big) \notag \\
   &= \Delta t \sum_{i} w_i\, g_i
        + \tfrac{1}{2}(\Delta t)^2 \sum_{i} w_i\, g_i^2
        \notag \\
   &\quad - \tfrac{1}{2}(\Delta t)^2
        \!\Big(\sum_{i} w_i\, g_i\Big)^{\!2}
        + O\!\big((g\,\Delta t)^3\big) \notag \\
   &= \Delta t \sum_{i} w_i\, g_i
        + \tfrac{1}{2}(\Delta t)^2 \!\left[\sum_{i} w_i\, g_i^2
        \right. \notag \\
   &\hspace{8em}\left.{}- \Big(\sum_{i} w_i\, g_i\Big)^{\!2}\right]
        + O\!\big((g\,\Delta t)^3\big).
  \label{eq:rp-taylor}
\end{align}
Dividing through by $\Delta t$ extracts the portfolio's annualized
one-step log-growth rate $g_p(t) := \Delta t^{-1}\log(W_t/W_{t-1})$,
the portfolio counterpart of the per-asset $g_i(t)$ (still
suppressing the time argument):
\begin{equation}
  g_p \;=\; \sum_{i} w_i\, g_i
   \;+\; \tfrac{1}{2}\,\Delta t \!\left(\sum_{i} w_i\, g_i^2
        \;-\; \Big(\sum_{i} w_i\, g_i\Big)^{\!2}\right)
   \;+\; O\!\big((\Delta t)^2\big).
  \label{eq:gp-onestep}
\end{equation}

To obtain expected annualized portfolio growth, we restored the time
argument and took unconditional expectations of
Eq.~\eqref{eq:gp-onestep} under stationarity. Define the mean and
variance of each asset's annualized growth rate:
\begin{equation}
  \mu_i \;:=\; \E\!\big(g_i(t)\big), \qquad
  \sigma_i^2 \;:=\; \Var\!\big(g_i(t)\big),
\end{equation}
the covariance of the annualized rates
$\Sigma_{ij} := \Cov\!\big(g_i(t), g_j(t)\big)$, and the portfolio
variance
$\sigma_p^2 := \Var\!\big(\textstyle\sum_i w_i g_i(t)\big) =
w^\top \Sigma w$. The $O(\Delta t)$ correction in
Eq.~\eqref{eq:gp-onestep} contains two squared quantities,
$\sum_i w_i\, g_i(t)^2$ and $(\sum_i w_i\, g_i(t))^2$. Both expand
by the second-moment identity, which holds for any random variable
$X$ with finite second moment:
\begin{equation}
  \E\!\big(X^2\big) \;=\; \E\!\big(X\big)^2 \;+\; \Var(X).
  \label{eq:second-moment}
\end{equation}
Applying this identity to $X = g_i(t)$ and summing against $w_i$ gives:
\begin{align}
  \E\!\Big(\sum_{i} w_i\, g_i(t)^2\Big)
   &\;=\; \sum_{i} w_i\, \E\!\big(g_i(t)^2\big)
       && \text{(linearity)} \notag \\
   &\;=\; \sum_{i} w_i\!\left[\,\mu_i^2 + \sigma_i^2\,\right]
       && \text{(by Eq.~\eqref{eq:second-moment})} \notag \\
   &\;=\; \sum_{i} w_i\, \mu_i^2 \;+\; \sum_{i} w_i\, \sigma_i^2.
   \label{eq:exp-gi2}
\end{align}
Applying the identity to $X = \sum_i w_i g_i(t)$, whose mean is
$\sum_i w_i\mu_i$ and whose variance is $\sigma_p^2$, gives:
\begin{align}
  \E\!\Big(\Big(\textstyle\sum_{i} w_i g_i(t)\Big)^{\!2}\Big)
   &\;=\; \E\!\Big(\textstyle\sum_{i} w_i g_i(t)\Big)^{\!2}
       \;+\; \Var\!\Big(\textstyle\sum_{i} w_i g_i(t)\Big)
       && \text{(by Eq.~\eqref{eq:second-moment})} \notag \\
   &\;=\; \Big(\textstyle\sum_{i} w_i\, \mu_i\Big)^{\!2}
       \;+\; \sigma_p^2.
   \label{eq:exp-sum2}
\end{align}
We substituted Eqs.~\eqref{eq:exp-gi2}--\eqref{eq:exp-sum2} into the
expectation of Eq.~\eqref{eq:gp-onestep} and grouped the variance and
mean contributions as follows:
\begin{equation}
\begin{aligned}
  \E(g_p) \;={}& \sum_{i} w_i\, \mu_i
  + \underbrace{\tfrac{1}{2}\,\Delta t\!\left(\sum_{i} w_i\, \sigma_i^2
       - \sigma_p^2\right)}_{\displaystyle D} \\
  &+ \tfrac{1}{2}\,\Delta t\!\left[\sum_{i} w_i\, \mu_i^2
       - \Big(\sum_{i} w_i\, \mu_i\Big)^{\!2}\right]
  + O\!\big((\Delta t)^2\big),
\end{aligned}
  \label{eq:booth-fama}
\end{equation}
the Booth--Fama decomposition, to second order, of the growth of a
daily-rebalanced portfolio in annualized growth-rate units
\cite{boothFama1992, markowitz1976longRun}. The middle term, $D$, is
the diversification return: $\tfrac{1}{2}\Delta t$ times the
difference between the weighted-average per-asset variance
$\sum_i w_i \sigma_i^2$ and the portfolio variance $\sigma_p^2$. The
moments $\mu_i$ and $\sigma_i^2$ of the annualized rates
have units $\mathrm{yr}^{-1}$ and $\mathrm{yr}^{-2}$ respectively;
the products $\Delta t\,\sigma_i^2$ and $\Delta t\,\sigma_p^2$ in
Eq.~\eqref{eq:booth-fama} are the annualized variances
($\mathrm{yr}^{-1}$). The third term in Eq.~\eqref{eq:booth-fama} is
the cross-sectional dispersion of the per-asset means, of order
$\Delta t\,\mu^2$, while $D$ is of order $\Delta t\,\sigma^2$. For
SPY, $\mu^2$ lies more than two orders of magnitude below $\sigma^2$
(Table~\ref{tab:descriptive_stats}), and we drop the third term,
treating the same ordering as typical of daily equity growth rates.
The realized annualized compound growth rate is
$\bar g_p = (T\,\Delta t)^{-1}\log(W_T/W_0) = T^{-1}\sum_{t} g_p(t)$.
If the composed growth rates form a stationary ergodic process,
$\bar g_p$ converges to $\E(g_p)$ by the ergodic theorem. The long-run
limit is then fixed by the unconditional joint distribution of
$(g_i(t))_i$ at a single date and does not depend on the serial
structure of the residual draw. This limit argument treats the
residuals as a stationary process; the per-path centering and scaling
of the composition act on a finite horizon and are outside it.

To show that $D$ is non-negative, we rewrote the bracket in $D$ as a
sum over asset pairs. Expanding the portfolio variance, using
$1 - w_i = \sum_{j \ne i} w_j$ for weights summing to one, and
collecting each pair once gives:
\begin{align}
  \sum_i w_i \sigma_i^2 - \sigma_p^2
   &\;=\; \sum_i w_i \sigma_i^2 - \sum_i w_i^2 \sigma_i^2
     - 2 \sum_{i<j} w_i w_j \Sigma_{ij}
     && \text{(expand $\sigma_p^2$)} \notag \\
   &\;=\; \sum_i w_i (1 - w_i)\, \sigma_i^2
     - 2 \sum_{i<j} w_i w_j \Sigma_{ij}
     && \text{(collect)} \notag \\
   &\;=\; \sum_i \sum_{j \ne i} w_i w_j\, \sigma_i^2
     - 2 \sum_{i<j} w_i w_j \Sigma_{ij}
     && \text{($1 - w_i = \textstyle\sum_{j \ne i} w_j$)} \notag \\
   &\;=\; \sum_{i<j} w_i w_j
     \left(\sigma_i^2 + \sigma_j^2 - 2\Sigma_{ij}\right)
     && \text{(pair terms)} \notag \\
   &\;=\; \sum_{i<j} w_i w_j \,\Var(g_i - g_j),
     && \text{(variance of a difference)}
  \label{eq:pair-identity}
\end{align}
so that:
\begin{equation}
  D \;=\; \tfrac{1}{2}\,\Delta t \sum_{i<j} w_i w_j \,\Var(g_i - g_j)
  \;\ge\; 0.
  \label{eq:dr-pairwise}
\end{equation}
The diversification return is therefore $\tfrac{1}{2}\Delta t$ times
a sum of pairwise difference variances with pair weights $w_i w_j$.
It is zero only when every pair of assets with positive weight
comoves perfectly ($\Var(g_i - g_j) = 0$) or all weight is on one
asset. The rebalancing rule itself produces
this effect: weights drift between rebalances, the rebalancer sells
whatever has outperformed (its weight rose above target) and buys
whatever has underperformed (its weight fell below target), and the
trade systematically converts period-by-period cross-sectional
dispersion into compounded growth. Markowitz
\cite{markowitz1976longRun} described this as the variance bonus of
the long-run mean-variance allocation; Booth and Fama
\cite{boothFama1992} called it the diversification return.

To separate market and residual contributions to the diversification
return, we substituted the single-index decomposition
$g_i(t) = \alpha_i + \beta_i\, g_m(t) + \eps_i(t)$ into
Eq.~\eqref{eq:dr-pairwise}, with market variance
$\sigma_m^2 := \Var(g_m(t))$, per-asset residual variance
$\sigma_{\eps,i}^2 := \Var(\eps_i(t))$, and $\Cov(g_m, \eps_i) = 0$
(checked empirically in Online Appendix~\ref{sec:supp-cov}). The
pairwise difference variance becomes:
\begin{equation}
  \Var(g_i - g_j) \;=\; (\beta_i - \beta_j)^2 \sigma_m^2
  \;+\; \Var(\eps_i - \eps_j).
\end{equation}
Because $g_m$ and the residuals are uncorrelated, this relation splits
the diversification return into market-beta-dispersion and
idiosyncratic components:
\begin{equation}
  D \;=\;
  \underbrace{\tfrac{1}{2}\,\Delta t\,\sigma_m^2 \sum_{i<j} w_i w_j (\beta_i - \beta_j)^2}_{D_{\mathrm{mkt}}}
  \;+\;
  \underbrace{\tfrac{1}{2}\,\Delta t \sum_{i<j} w_i w_j \,\Var(\eps_i - \eps_j)}_{D_\eps},
  \label{eq:dr-sim}
\end{equation}
where the market-beta-dispersion component $D_\mathrm{mkt}$ contains
the weighted cross-sectional variance of $\beta$. Applying the pair
identity of Eq.~\eqref{eq:pair-identity} with $\beta_i$ in place of
$g_i$ gives:
\begin{align}
  \sum_{i<j} w_i w_j (\beta_i - \beta_j)^2
   &\;=\; \sum_i w_i (1 - w_i)\, \beta_i^2
     - 2 \sum_{i<j} w_i w_j \beta_i \beta_j
     && \text{(pair identity)} \notag \\
   &\;=\; \sum_i w_i \beta_i^2
     - \Big(\sum_i w_i^2 \beta_i^2
     + 2 \sum_{i<j} w_i w_j \beta_i \beta_j\Big)
     && \text{(regroup)} \notag \\
   &\;=\; \sum_i w_i \beta_i^2 - \Big(\sum_i w_i \beta_i\Big)^{\!2},
     && \text{(square of a sum)}
\end{align}
so $D_\mathrm{mkt} = \tfrac{1}{2}\,\Delta t\,\sigma_m^2 [\sum_i w_i
\beta_i^2 - (\sum_i w_i \beta_i)^2]$, the weighted variance of $\beta$
scaled by the annualized market variance $\Delta t\,\sigma_m^2$.
Across the $423$ non-market tickers, the calibrated $\beta$
distribution spans $[0.03, 1.79]$ (Fig.~\ref{fig:branch_map}). A
weighted variance is at most one quarter of the squared range, so
$D_{\mathrm{mkt}} \le \tfrac{1}{8}\,\Delta t\,\sigma_m^2\,(1.79 -
0.03)^2$. The SPY closing-price series used for the multi-asset
universe has annualized variance $\Delta t\,\sigma_m^2 =
0.030~\mathrm{yr}^{-1}$ (Fig.~\ref{fig:preservation}), at which
$D_{\mathrm{mkt}}$ is at most $0.012~\mathrm{yr}^{-1}$, or $1.2$
percentage points per year. For any long-only weights on this
universe, the market component therefore contributes little, and a
larger diversification return must come from the idiosyncratic
component $D_\eps$. In the synthetic construction, residual paths are
generated independently across assets, so
$\Var(\eps_i - \eps_j) = \sigma_{\eps,i}^2 + \sigma_{\eps,j}^2$ and,
by the same pair identity,
$D_\eps = \tfrac{1}{2}\,\Delta t \sum_i w_i (1 - w_i)\,
\sigma_{\eps,i}^2$, a weighted sum of the annualized idiosyncratic
variances $\Delta t\,\sigma_{\eps,i}^2$.

Because the second-order value of $D$ depends only on same-date
unconditional moments, serial dependence in the residual draw does not
change it. Equation~\eqref{eq:scale} chooses $s_i$ so that each
composed variance $\beta_i^2\sigma_m^2+\sigma_{\eps,i}^2$ equals the
generator variance $\sigma_{\gen,i}^2$, and any residual generator
with the same same-date variances gives the same $D$. As an
illustration, consider a stationary first-order autoregressive
[AR(1)] residual $\eps_i(t) = \phi\, \eps_i(t-1) + \eta_i(t)$ with
independent and identically distributed innovations $\eta_i$ whose
variance is chosen so that
$\Var(\eps_i) = \sigma_{\eta,i}^2 / (1 - \phi^2) = \sigma_{\eps,i}^2$.
The parameter $\phi$ leaves the variance at a single date unchanged,
so it does not change $\E(g_p)$ in Eqs.~\eqref{eq:booth-fama}
and~\eqref{eq:dr-sim}. It does change sampling noise: for an
uncentered stationary AR(1) residual with variance $\sigma_\eps^2$,
the variance of the sample mean is approximately
$T^{-1}\sigma_\eps^2(1+\phi)/(1-\phi)$. Serial dependence therefore
changes the per-path noise around the stationary mean, not the mean
itself, and is not a lever on the diversification return at the level
of the second-order formula.

The difference between synthetic and real residuals that does change
$D$ is contemporaneous. The derivation assumes residuals are
uncorrelated across assets at each date, as they are in the
independent per-asset generator. Real residuals share sector and style
factors, which reduce $\Var(\eps_i-\eps_j)$ below the
independent-generator value $\sigma_{\eps,i}^2+\sigma_{\eps,j}^2$.
Matching each marginal variance does not restore this covariance, so
the synthetic $D_\eps$ exceeds its real-market counterpart when the
omitted covariance is positive. We did not measure the size of this
excess, which would require a joint-residual experiment with the
marginals and the allocation rule held fixed. A comparison with
realized portfolio returns would also contain a second effect:
realized diversification returns on real equity portfolios fall short
of the second-order formula at all rebalancing
horizons~\cite{boucheyNemtchinovWong2015}. A measured excess of
synthetic over realized diversification return would therefore bound
the part caused by independent residuals from above rather than
measure it.


% =============================================================================
% S3 Validation metric definitions (KS, AD, W1, Hellinger)
% =============================================================================

\section{Validation metric definitions}
\label{sec:supp-validation}

To distinguish marginal agreement from temporal agreement, we used
separate distributional, tail, and autocorrelation summaries
(Tables~\ref{tab:model_comparison} and~\ref{tab:jump_ablation}).
This section defines those summaries; the factor-recovery errors and
variance ratios of the composition scorecard are defined in the main
text with that experiment. For one observed series $x=(x_1,\ldots,x_n)$ and $R$ synthetic
paths $y^{(r)}$ of the same length, each distributional comparison
used one complete synthetic path and the fixed observed history.
The two-sample Kolmogorov--Smirnov (KS) statistic is the largest
absolute difference between their empirical cumulative distribution
functions, $D_{\mathrm{KS}}=\sup_z|F_x(z)-F_{y^{(r)}}(z)|$
\cite{kolmogorov1933,smirnov1948}. The Anderson--Darling (AD)
comparison gives greater weight to differences near the tails
\cite{andersonDarling1952}. We used the approximate two-sample KS
test and the modified, tie-aware two-sample AD calculation in
\texttt{HypothesisTests.jl}, with asymptotic $p$-values. A path
passed when $p>0.05$, and the pass percentage was $100$ times the
fraction of passing paths. These percentages are not calibrated
significance levels, for three reasons: returns are serially
dependent, whereas the tests assume independent observations; every
path is compared with the same observed series, so the comparisons
are not independent of one another; and in sample, the generator was
fitted to that series. A high pass rate therefore describes agreement
with the supplied history and is not evidence that the model and the
data share a distribution.

To quantify marginal discrepancies without a pass threshold, we
computed Wasserstein-1 and histogram Hellinger scores for each path
(Table~\ref{tab:model_comparison}). Let $x_{(j)}$ and $y^{(r)}_{(j)}$
denote the $j$th order statistics in increasing order. Because each
synthetic path has the same length as the observed series, the
Wasserstein-1 distance is:
\begin{equation}
 W_1(x,y^{(r)})=\frac{1}{n}\sum_{j=1}^{n}
       |x_{(j)}-y^{(r)}_{(j)}|.
 \label{eq:wasserstein}
\end{equation}
It has the growth-rate units $\mathrm{yr}^{-1}$. For each Hellinger
comparison, we formed $50$ equally spaced bins between the minimum
and maximum of the pooled pair of samples. Let $p_b$ and $q_b$ be
the observed and simulated bin counts divided by the sample length
$n$. The histogram score is:
\begin{equation}
 H(x,y^{(r)})=\left[\frac{1}{2}\sum_{b=1}^{50}
          (\sqrt{p_b}-\sqrt{q_b})^2\right]^{1/2}.
 \label{eq:hellinger}
\end{equation}
Every bin, including the last, contained its lower edge and excluded
its upper edge, so observations equal to the pooled maximum fell
outside all bins and that sample's bin weights sum to slightly less
than one; we state this convention because the reported scores were
computed with it. The score is a histogram approximation to the
Hellinger distance between the two densities. Its bin edges depend on
each pair, and its bins hold fewer observations in shorter samples, so
it is best compared across generators within one window; holdout
scores on $249$ observations run larger than training scores partly
for that reason. Smaller Wasserstein and Hellinger scores indicate
closer marginal agreement.

To describe tail shape, we computed excess kurtosis for each
synthetic path and compared its mean over paths with the observed value
(Table~\ref{tab:model_comparison}); the composition tables report
medians instead. With
$m_p(x)=n^{-1}\sum_{t=1}^{n}(x_t-\bar x)^p$, the estimator was
$\kappa(x)=m_4(x)/m_2(x)^2-3$, without a finite-sample bias
correction. The composition scorecard also used the Hill index of
absolute returns (Table~\ref{tab:aggregate}). Let
$X_{(1)}\ge\cdots\ge X_{(n)}$ be decreasingly ordered absolute
returns and $k=\max\{2,\lfloor0.05n\rfloor\}$. The implemented
upper-tail estimator is:
\begin{equation}
 \hat\xi=\frac{1}{k-1}\sum_{j=1}^{k-1}
             \log\!\left(\frac{X_{(j)}}{X_{(k)}}\right).
 \label{eq:hill}
\end{equation}
The estimate $\hat\xi$ is the reciprocal of the conventional tail
index, so larger values indicate a heavier absolute-return tail; the
threshold $X_{(k)}$ must be positive. These shape summaries do not
measure how often returns exceed a Value-at-Risk threshold.

To assess whether the simulated distributions bracketed the observed
one, we measured agreement at $99$ quantile levels
(Table~\ref{tab:model_comparison}). For $u_j=j/100$, $j=1,\ldots,99$,
let $Q_x(u_j)$ be the observed sample quantile and
$Q_r(u_j)$ the corresponding quantile of path $r$. Let
$L_j$ and $U_j$ be the $5$th and $95$th percentiles of
$\{Q_r(u_j):r=1,\ldots,R\}$. All quantiles used linear interpolation
between order statistics, the common software default (type~7). The reported
quantile-envelope coverage is:
\begin{equation}
 C=\frac{100}{99}\sum_{j=1}^{99}
       \mathbf{1}\{L_j\le Q_x(u_j)\le U_j\}.
 \label{eq:quantile-envelope-coverage}
\end{equation}
Larger values mean that more observed quantiles fall inside the
pointwise simulation envelopes. This statistic assesses marginal
quantile agreement; it does not estimate daily Value-at-Risk
exceedance coverage or simultaneous confidence coverage.

To quantify persistence in return magnitude, we used the sample
autocorrelation function (ACF) of absolute returns. Each series was
centered by its own full-series absolute-return mean,
$\bar a_x=n^{-1}\sum_t|x_t|$. At lag $\ell$, the calculation was:
\begin{equation}
 a_\ell(x)=
 \frac{\sum_{t=1}^{n-\ell}(|x_t|-\bar a_x)(|x_{t+\ell}|-\bar a_x)}
      {\sum_{t=1}^{n}(|x_t|-\bar a_x)^2}.
 \label{eq:sample-absolute-acf}
\end{equation}
The denominator used all $n$ observations, with no lag-dependent
rescaling. Both error scores below start at lag one.
The single-asset comparison measured the error of the ensemble-average
ACF (Table~\ref{tab:model_comparison}). For a lag limit $L$, that
mean absolute error is:
\begin{equation}
 E^{\rm ensemble}_L=\frac{1}{L}\sum_{\ell=1}^{L}
 \left|a_\ell(x)-\frac{1}{R}\sum_{r=1}^{R}a_\ell(y^{(r)})\right|.
 \label{eq:acf-ensemble-error}
\end{equation}
We used $R=1{,}000$ and $L=\min(252,n-1)$: lags $1$--$252$
for the $2{,}766$-observation training series and $1$--$248$ for the
$249$-observation holdout. This evaluation lag limit is separate from
the $25$-lag window of the tuning objective in
Eq.~\eqref{eq:grid_search}. The state-resolution and emission
comparisons (Tables~\ref{tab:sensitivity_n}
and~\ref{tab:emission_hsmm}) used the same ensemble convention.

The multi-asset jump comparison instead measured error within each
path before averaging (Table~\ref{tab:jump_ablation}). For one
asset, its mean path error is:
\begin{equation}
 E^{\rm path}_L=\frac{1}{R}\sum_{r=1}^{R}
       \left[\frac{1}{L}\sum_{\ell=1}^{L}
       |a_\ell(x)-a_\ell(y^{(r)})|\right].
 \label{eq:acf-path-error}
\end{equation}
We used $L=25$ and $60$ in both windows and averaged these scores
over assets. The fallback sensitivity check also reported this path
score with $L=25$ (Table~\ref{tab:fallback_sensitivity}). At a common
lag limit and on the same paths, the triangle inequality gives
$E^{\rm ensemble}_L\le E^{\rm path}_L$: averaging the curves can
cancel errors of opposite sign. The ensemble score asks whether the
model's average autocorrelation curve matches the observed curve; the
path score asks how far a single path's curve typically falls from
it, sampling noise included. Their magnitudes
cannot be compared directly across the single-asset and composition
experiments, which also differ in lag limits, histories, and assets.

To describe simulation uncertainty in the single-asset comparison,
we used a standard-error method suited to each statistic
(Table~\ref{tab:model_comparison}). For a pass fraction $\hat p$,
the reported binomial standard error was
$100\sqrt{\hat p(1-\hat p)/R}$ percentage points. This estimate is
zero when every path passes or fails; it does not establish that the
underlying pass probability is exactly zero or one. For kurtosis,
Wasserstein distance, and Hellinger score, the standard error was
the sample standard deviation of the $R$ path scores divided by
$\sqrt R$. For ensemble ACF error and quantile-envelope coverage,
we used $500$ bootstrap resamples of the $R$ path indices, each of
size $R$ with replacement. Each resample retained entire paths and
recomputed the statistic across all lags or quantile levels. The
standard error was the sample standard deviation of those $500$
recomputed values, without a further division by $\sqrt{500}$.
All these errors condition on one observed history and a fixed fitted
model; they omit uncertainty in fitting, model selection, and market
regimes. The multi-asset jump contrasts used the replication as the
unit instead: within each of the $1{,}000$ replications we averaged a
metric across assets, differenced two jump settings, and took the
standard deviation of the $1{,}000$ differences divided by
$\sqrt{1{,}000}$ as the standard error, so that the dependence created
by the shared market path stays inside each unit
(Table~\ref{tab:jump_ablation_uncertainty}; Online
Appendix~\ref{sec:supp-jump-composition}).

% --- Figure: Statistical validation of HMM-WJ ---
\begin{figure}[H]
  \centering
  \includegraphics[width=\textwidth]{figs/supplement/FigS04-Statistical-Validation.pdf}
  \caption{\textbf{Statistical validation of the hidden Markov model with
    jumps (HMM-WJ)} ($N = 100$,
    $\alpha = 0.05$). Panels (a) and (d) condition on the
    $\sim\!24\%$ of in-sample paths that contain at least one jump
    episode. (a)~In-sample
    KS $p$-values for these $239$ paths (pass rate $93.7\%$). (b)~Out-of-sample
    KS $p$-values for all $1{,}000$ paths (pass rate $95.4\%$), since few
    holdout-length paths contain an episode. (c)~Out-of-sample
    density fan chart; the observed density falls within the simulation
    envelope. (d)~Observed and mean simulated in-sample autocorrelation
    function (ACF) of $|G_t|$
    for the paths with an episode.}
  \label{fig:statistical_validation}
\end{figure}


\subsection{Six-method composition evaluation}
\label{sec:supp-composition-evaluation}

The four residual-based comparators used the observed asset's
ordinary least-squares (OLS) residuals to set their distributions
(Table~\ref{tab:aggregate}). The Gaussian single-index model (SIM)
used independent Gaussian draws with the OLS residual variance. The
residual-fit hidden Markov model, labeled JumpHMM-on-residuals after
the software package, fitted a no-jump generator to a synthetic price
path obtained by cumulative compounding of the residuals. The
Politis--Romano stationary bootstrap~\cite{politisRomano1994} used a
mean block length of $50$ days in both training and holdout. The
generalized autoregressive conditional heteroskedasticity model,
GARCH(1,1)-$t$, fitted conditional variance to the residual
series~\cite{bollerslev1986}. We accepted a fit only if optimization
converged, the likelihood and coefficients were finite, the variance
intercept was positive, and the nonnegative ARCH and GARCH
coefficients summed to less than one. The Student-$t$ degrees of
freedom had to exceed two, the unconditional variance had to be finite
and positive, and a trial simulated path had to be finite and
nonconstant. All $423$ optimizations converged, but $30$ failed the
coefficient-sum condition, leaving $393$ eligible assets; stored fits
were rechecked before evaluation. All four residual
comparators entered composition without further centering or
rescaling.

The six-method scores measured agreement with the observed asset and
its calibrated factor model, rather than recovery of a common target
for every branch (Tables~\ref{tab:aggregate}, \ref{tab:branch}, and
\ref{tab:beta_bucket}). Lower absolute OLS errors indicated closer
recovery, but on the hybrid variance-preserving branch the composed
$R^2$ approached the draw's $\rho_i$, capped at $1-f$, rather than
$R^2_{i,\real}$. Its $R^2$ error therefore mainly measured this
designed shift; loading error also included any change made by
clipping. A composed-to-generator variance ratio of one indicated
equal variances, which the $R^2$-preserving branch did not target.
Higher KS and AD pass rates and lower Wasserstein-1 distance indicated
closer distributional agreement; kurtosis and Hill indices were
compared with the observed asset's values. We omitted cross-sectional
rank reordering to test each ticker's marginal construction and
scored volatility clustering separately in the jump-enabled comparison.
The paired centering control isolated differences in residual-mean
constraints (Online Appendix~\ref{sec:supp-centering-control}).


\subsection{Gated recurrent unit benchmark}
\label{sec:supp-gru}

To provide an autoregressive neural benchmark, we used a gated
recurrent unit (GRU) network to predict the distribution of the next
SPY growth rate from the preceding $50$ values
(Table~\ref{tab:model_comparison}). We standardized the training
series using its mean and population standard deviation and reused
these two values unchanged to standardize inputs and to convert
generated values back to growth rates. The network had two GRU layers
with $64$ hidden units per layer, dropout probability $0.1$ between
layers during training, and separate linear heads for a conditional
mean and log variance. Both heads read the last hidden output of each
input window, and the network had $37{,}954$ trainable parameters.
For the standardized target $z_t$ that follows a window, with
predicted mean $m_t$ and predicted log variance $v_t$, the training
loss was the Gaussian negative log-likelihood with its constant term
omitted:
\begin{equation}
 \mathcal L=\frac{1}{2B}\sum_{t\in\mathcal B}
       \left[v_t+(z_t-m_t)^2\exp(-v_t)\right],
 \label{eq:gru-loss}
\end{equation}
where $\mathcal B$ is a training batch containing $B=256$ windows.
The $2{,}766$ training observations supplied $2{,}716$ overlapping
input-target pairs. Each epoch shuffled these pairs and used ten
complete batches; the final incomplete batch was dropped. The Adam
optimizer used an initial learning rate of $10^{-3}$, default moment
decay rates $(0.9,0.999)$, and no weight decay. We clipped the
gradient norm at one and halved the learning rate whenever mean
training loss had not improved for $15$ epochs, down to a minimum
learning rate of $10^{-5}$. Training ran for at most $200$ epochs. We
retained weights whenever mean training loss improved by more than
$10^{-4}$ and would have stopped after $40$ epochs without such
improvement. That early-stopping rule never triggered: the fit ran
the full $200$ epochs, and the best training loss occurred at epoch
$199$. Model selection used training loss, with no separate
validation split; the $2025$ holdout played no part in fitting or
selection.

To generate each path, we disabled dropout and drew its initial
window uniformly from the $2{,}716$ observed training input windows.
At each step, we passed the current $50$ values through the network
with a fresh zero hidden state, sampled a Gaussian value with mean
$m_t$ and standard deviation $\max\{\exp(v_t/2),10^{-6}\}$, and
rolled the window forward, dropping its oldest value and appending
the draw. A path consisted of the generated values only, with no
burn-in, and the fixed training normalization was inverted afterward.
The benchmark was implemented in Python. The same fitted network
generated $1{,}000$ training-length and then $1{,}000$ holdout-length
paths, in that order, from one random stream seeded with $1234$ for
NumPy and PyTorch, so the order is part of the reproduction recipe.
To confirm reproducibility, we retrained the network and regenerated
both path sets on a central processing unit (CPU) with Python 3.14.2,
PyTorch 2.10.0, NumPy 2.4.3, and SciPy 1.17.1. This rerun generated
paths in parallel, starting each worker from the random-number state
that one-at-a-time generation would have reached at that path, and it
reproduced the original training-length and holdout-length path
files exactly. The repository includes pinned dependencies, trained
weights, an epoch log, and input and output fingerprints. The
benchmark is one fitted network, so its reported Monte Carlo standard
errors do not measure variation across training seeds or neural
architectures.

% =============================================================================
% S4  Derivation of the single-asset generator
% Six subsections: generative law; estimators; stationary law and moments;
% dependence without jumps; jump-episode arithmetic; what (epsilon, lambda)
% control. Every numerical value evaluates the stated closed form at the
% fitted SPY model (generator script: code/spy-experiment/SI-Derivation-Checks.jl).
% =============================================================================

\section{Derivation of the single-asset generator}
\label{sec:supp-generator}

This appendix derives the probability law of the fitted single-asset
generator and connects each fitted component to the empirical feature
it controls. The hidden Markov model with jumps (HMM-WJ), its no-jump
restriction (HMM-NJ), and their symbols follow the main-text
definitions in Eq.~\eqref{eq:emission} and
Algorithm~\ref{alg:hmmwj}. Unless noted otherwise, numerical values
evaluate the closed-form expressions at the fitted SPY model of
Table~\ref{tab:hyperparams}: $N = 100$ states, Student-$t$ degrees of
freedom $\nu = 5$, jump probability $\epsilon = 10^{-4}$, mean episode
duration $\lambda = 90$ steps, tail-selection probability
$p_{\rm neg} = 0.52$, tail-set size $N_{\rm tail} = 5$, and the
in-sample series of $T = 2{,}766$ daily excess growth rates described
in Table~\ref{tab:descriptive_stats}.
Table~\ref{tab:supp-generator-quantities} collects the closed-form
quantities evaluated at these values beside their sample or simulated
counterparts.

% --- Table: SPY descriptive statistics ---
\begin{table}[H]
\centering
\caption{\textbf{Descriptive statistics and stylized-facts tests for
    SPY daily excess growth rates.} The series is computed from daily
    volume-weighted average prices, the SPY series used to fit the
    single-asset generator. In-sample (IS): $2014$--$2024$
    ($T = 2{,}766$); out-of-sample (OoS): $2025$ ($T = 249$). The
    growth rate $G_t$ is the daily log excess return divided by
    $\Delta t = 1/252$~yr, so its standard deviation is not an
    annualized volatility; the annualized volatility is
    $\sigma(G_t)\sqrt{\Delta t}$, about $13.5\%$ in sample.
    Gaussian and Laplace columns show maximum likelihood estimation
    (MLE) fits to the in-sample data. The Laplace MLE location is the
    sample median and its scale is the mean absolute deviation about
    the median (Eq.~\eqref{eq:supp-laplace-mle}), so its mean and
    standard deviation differ from the sample moments. We applied the
    Jarque--Bera (JB) normality test and the Ljung--Box (LB)
    autocorrelation test at lag $20$ to the observed series only.
    $^{\ast}$ denotes rejection at $\alpha = 0.05$.}
\label{tab:descriptive_stats}
\small
\begin{tabular}{@{}lcccc@{}}
\toprule
\textbf{Statistic} & \textbf{IS Observed} & \textbf{OoS Observed} & \textbf{Gaussian} & \textbf{Laplace} \\ \midrule
Mean (annualized, \%)       & 6.31   & 11.60  & 6.31   & 14.19  \\
Standard deviation of $G_t$ (\% yr$^{-1}$) & 214.50 & 220.62 & 214.46 & 205.93 \\
Skewness                    & $-0.753$ & $-1.093$ & 0    & 0      \\
Excess Kurtosis             & 7.715  & 6.867  & 0      & 3      \\
JB normality test           & reject$^\ast$ ($p<0.001$) & reject$^\ast$ ($p<0.001$) & not used & not used \\
LB test on $G_t$ (lag 20)   & reject$^\ast$ ($p<0.001$) & reject$^\ast$ ($p=0.019$) & not used & not used \\
LB test on $|G_t|$ (lag 20) & reject$^\ast$ ($p<0.001$) & reject$^\ast$ ($p<0.001$) & not used & not used \\
\bottomrule
\end{tabular}
\end{table}

% --- Table: closed-form quantities at the fitted SPY generator ---
\begin{table}[H]
\centering
\caption{\textbf{Closed-form quantities of the fitted SPY generator
    against sample and simulated values.} Each closed-form entry
    evaluates an expression derived in this section at the fitted SPY
    model of Table~\ref{tab:hyperparams}. Sample values are computed
    from the in-sample SPY series; simulated values use $400$
    in-sample-length paths for the jump-episode quantities (forced-step
    fraction, jump-path probability, and forced fraction within a jump
    path) and $1{,}000$ paths for skewness. The sup-distance entry is
    the largest gap between the model's stationary distribution
    function and the in-sample empirical distribution function; for
    comparison, the two-sample Kolmogorov--Smirnov critical value used
    by the path-level tests at $\alpha = 0.05$ for two samples of
    $T = 2{,}766$ observations is $1.358\sqrt{2/T} = 0.0365$. Observed
    autocorrelations are those of the in-sample series shown in
    Fig.~\ref{fig:empirical_motivation}d.}
\label{tab:supp-generator-quantities}
\footnotesize
\resizebox{\textwidth}{!}{\input{sections/tables/tableS_generator_quantities}}
\end{table}

\subsection{Generative law}
\label{sec:supp-gen-law}

The no-jump model generates a hidden state path $S_0, \dots, S_T$ on
$\{1, \dots, N\}$ and observations $G_1, \dots, G_T$ with the joint
density:
\begin{equation}
  p(s_0, \dots, s_T,\, g_1, \dots, g_T)
  \;=\; \bar\pi_{s_0} \prod_{t=1}^{T}
  \mathbf{T}_{s_{t-1},\, s_t}\; f_{s_t}(g_t),
  \label{eq:supp-joint-law}
\end{equation}
where $\bar\pi_{s_0}$ is the initial probability of state $s_0$,
$\mathbf{T}_{s_{t-1}, s_t}$ is the probability of moving from state
$s_{t-1}$ to state $s_t$, and $f_k$ is the emission density of state
$k$; bold $\mathbf{T}$ is the transition matrix and italic $T$ the
number of observations. Because $\bar\pi$ is stationary for
$\mathbf{T}$, starting from an unemitted $S_0 \sim \bar\pi$ gives
$S_1 \sim \bar\pi$, the same law as drawing $S_1$ from $\bar\pi$ and
emitting it, as Algorithm~\ref{alg:hmmwj} does. Under
Eq.~\eqref{eq:emission}, $f_k$ is the location-scale
Student-$t$ density:
\begin{equation}
  f_k(g) \;=\; \frac{1}{\sigma_k}\,
  \psi_\nu\!\left(\frac{g - \mu_k}{\sigma_k}\right),
  \qquad
  \psi_\nu(x) \;=\;
  \frac{\Gamma\!\big(\tfrac{\nu+1}{2}\big)}
       {\sqrt{\nu \pi}\; \Gamma\!\big(\tfrac{\nu}{2}\big)}
  \left(1 + \frac{x^2}{\nu}\right)^{\!-\frac{\nu+1}{2}},
  \label{eq:supp-emission-density}
\end{equation}
where $\psi_\nu$ is the density of a standard Student-$t$ random
variable with $\nu$ degrees of freedom, $\Gamma$ is the gamma
function, and $(\mu_k, \sigma_k)$ are the location and scale of state
$k$. The factorization in Eq.~\eqref{eq:supp-joint-law} states the
structural property that drives everything that follows: given the
state path, the observations are independent. Every form of serial
dependence the generator can produce must therefore be carried by the
state path~\cite{rabinerJuang1986}.

The jump-augmented model enlarges the hidden state. Let
$J_t \in \{0, 1, 2, \dots\}$ count the forced steps remaining after
time $t$, starting from $J_0 = 0$, and let $u$ be the forced-step
state distribution, which
mixes uniform draws over the two tail sets:
\begin{equation}
  u_k \;=\;
  \begin{cases}
    p_{\rm neg}/N_{\rm tail} & k \in \mathcal{S}_{-},\\[2pt]
    (1 - p_{\rm neg})/N_{\rm tail} & k \in \mathcal{S}_{+},\\[2pt]
    0 & \text{otherwise}.
  \end{cases}
  \label{eq:supp-u}
\end{equation}
Writing $p_K(m) = e^{-\lambda} \lambda^m / m!$ for the Poisson
probability of episode length $m$ and $\mathbf{1}\{\cdot\}$ for the
indicator function, the pair $(S_t, J_t)$ evolves by the transition
kernel:
\begin{equation}
\begin{aligned}
  &\Pr\big(S_t = k,\, J_t = j \mid S_{t-1} = s,\, J_{t-1} = i\big) \\
  &\qquad=\;
  \begin{cases}
    u_k\, \mathbf{1}\{j = i - 1\} & i \ge 1,\\[6pt]
    \big(1 - \epsilon + \epsilon e^{-\lambda}\big)\,
      \mathbf{T}_{s, k}\, \mathbf{1}\{j = 0\}
      \;+\; \epsilon\, p_K(j + 1)\, u_k & i = 0,
  \end{cases}
\end{aligned}
  \label{eq:supp-kernel}
\end{equation}
and the emission draw depends only on the first coordinate,
$G_t \sim f_{S_t}$. The two cases read as follows. While an episode is
active ($i \ge 1$), the next state is an independent draw from $u$ and
the counter decrements; the tail set and the state within it are
re-drawn at every forced step. While no episode is active ($i = 0$),
the chain either takes an ordinary transition (no trigger, or a
trigger whose Poisson draw returns zero) or starts an episode of
length $j + 1 \ge 1$ whose first forced step occurs immediately. The
pair $(S_t, J_t)$ is therefore itself a Markov chain, and the HMM-WJ
generator is a hidden Markov model on the enlarged state space
$\{1, \dots, N\} \times \{0, 1, 2, \dots\}$. Algorithm~\ref{alg:hmmwj}
implements Eq.~\eqref{eq:supp-kernel} with the episode counter
truncated at the simulation horizon. It emits the initial draw from
$\bar\pi$ as the first observation and applies the kernel from the
second step onward, so the first observed day can never be a forced
step; this is the only way its path law differs from the kernel's.
When an episode ends, ordinary transitions resume from the last forced
tail state, so the chain returns to the central states through the
tail rows of $\mathbf{T}$. That return is the only link between an
episode and the free steps that follow it, which the autocorrelation
derivation of Online Appendix~\ref{sec:supp-gen-tuning} uses.

\subsection{Estimators and the one-step empirical law}
\label{sec:supp-gen-estimators}

The fitting procedure consists of four closed-form estimators: the
Laplace partition, the counted transition matrix, the within-state
emission moments, and the initial distribution. Each is a function of
a simple sample summary, and each reproduces that summary exactly or
up to an end effect of order $1/T$. The Laplace partition comes
first. For observations $G_1, \dots, G_T$, the Laplace
log-likelihood with location $\mu_L$ and scale $b_L$ is:
\begin{equation}
  \ell(\mu_L, b_L) \;=\;
  -T \ln (2 b_L) \;-\; \frac{1}{b_L} \sum_{t=1}^{T} |G_t - \mu_L|.
  \label{eq:supp-laplace-ll}
\end{equation}
For any fixed $b_L > 0$, maximizing $\ell$ over $\mu_L$ minimizes
$\sum_t |G_t - \mu_L|$, so $\hat\mu_L$ is the sample median. Setting
the $b_L$ derivative to zero at $\hat\mu_L$ gives:
\begin{align}
  \frac{\partial \ell}{\partial b_L}
   &\;=\; -\frac{T}{b_L}
     \;+\; \frac{1}{b_L^2} \sum_{t=1}^{T} |G_t - \hat\mu_L|
     \;=\; 0
     && \text{(first-order condition)} \notag \\
  \hat b_L
   &\;=\; \frac{1}{T} \sum_{t=1}^{T} |G_t - \hat\mu_L|,
     && \text{(solve for $b_L$)}
  \label{eq:supp-laplace-mle}
\end{align}
the mean absolute deviation about the median. The state boundaries
then invert the fitted Laplace cumulative distribution function
(CDF)~\cite{kotzKozubowskiPodgorski2001}. The Laplace CDF is
piecewise exponential:
\begin{equation}
  F_L(x) \;=\;
  \begin{cases}
    \tfrac{1}{2} \exp\!\big((x - \hat\mu_L)/\hat b_L\big)
      & x \le \hat\mu_L,\\[4pt]
    1 - \tfrac{1}{2} \exp\!\big(-(x - \hat\mu_L)/\hat b_L\big)
      & x > \hat\mu_L,
  \end{cases}
  \label{eq:supp-laplace-cdf}
\end{equation}
and setting $F_L(Q_k) = k/N$ and solving the lower branch for
$k \le N/2$ and the upper branch for $k \ge N/2$ gives the closed-form
boundaries:
\begin{align}
  \frac{k}{N} &\;=\; \tfrac{1}{2}\,
    e^{(Q_k - \hat\mu_L)/\hat b_L}
  \;\Longrightarrow\;
  Q_k \;=\; \hat\mu_L + \hat b_L \ln\!\frac{2k}{N},
    && \text{(lower branch)} \notag \\
  \frac{k}{N} &\;=\; 1 - \tfrac{1}{2}\,
    e^{-(Q_k - \hat\mu_L)/\hat b_L}
  \;\Longrightarrow\;
  Q_k \;=\; \hat\mu_L - \hat b_L
    \ln\!\Big(2 - \tfrac{2k}{N}\Big).
    && \text{(upper branch)}
  \label{eq:supp-quantiles}
\end{align}
The states are equiprobable under the fitted Laplace, not under the
empirical distribution, so the in-sample state counts are not uniform:
across the $100$ SPY states, they range from $16$ to $47$ of the
$2{,}766$ observations (Table~\ref{tab:supp-generator-quantities}).
For SPY, every state therefore held at least $16$ observations with a
nonzero spread, and neither fallback rule for the emission parameters
applied.

The transition matrix is the maximum likelihood estimator (MLE) of a
first-order Markov chain. Conditional on the assigned state sequence
$s_1, \dots, s_T$, let $n_{jk}$ count the transitions from state $j$
to state $k$ and let $n_{j \cdot} = \sum_k n_{jk}$ be the row total.
The likelihood of the transition parameters factorizes over rows into
multinomials, and maximizing the log-likelihood subject to each row
summing to one, with multiplier $c_j$ for row $j$, uses the
Lagrangian:
\begin{equation}
  \mathcal{L} \;=\; \sum_{j, k} n_{jk} \ln \mathbf{T}_{jk}
    \;+\; \sum_{j} c_{j} \Big(1 - \sum_{k} \mathbf{T}_{jk}\Big),
  \label{eq:supp-lagrangian}
\end{equation}
whose first-order conditions solve as follows:
\begin{align}
  \frac{\partial \mathcal{L}}{\partial \mathbf{T}_{jk}}
   &\;=\; \frac{n_{jk}}{\mathbf{T}_{jk}} - c_j \;=\; 0
     && \text{(first-order condition)} \notag \\
  \mathbf{T}_{jk}
   &\;=\; \frac{n_{jk}}{c_j}
     && \text{(solve for $\mathbf{T}_{jk}$)} \notag \\
  c_j
   &\;=\; \sum_{k} n_{jk} \;=\; n_{j \cdot}
     && \text{(row constraint $\textstyle\sum_k \mathbf{T}_{jk} = 1$)} \notag \\
  \widehat{\mathbf{T}}_{jk}
   &\;=\; \frac{n_{jk}}{n_{j \cdot}}.
     && \text{(counted estimator)}
  \label{eq:supp-markov-mle}
\end{align}
For a row with $n_{j\cdot}>0$, the first-order condition applies to
cells with positive counts. A cell with $n_{jk}=0$ adds nothing to the
log-likelihood, so any probability placed on it is taken from observed
transitions and lowers the likelihood; its maximizer is therefore
$\widehat{\mathbf{T}}_{jk}=0$, which the same formula returns.
When $n_{j\cdot}=0$, the observations do not identify that row.
The implementation sets every entry to $1/N$; this is a fallback
choice, not an estimate of an observed transition frequency.
Multiplying $\widehat{\mathbf{T}}_{jk}$ by the empirical origin
frequency $n_{j\cdot}/(T-1)$ recovers $n_{jk}/(T-1)$, the empirical
frequency of the state pair $(j, k)$. Paired with the empirical
origin frequencies, the counted estimator therefore reproduces the
in-sample joint law of $(S_{t-1}, S_t)$ exactly; everything it
asserts about longer horizons is the Chapman--Kolmogorov extension
$\mathbf{T}^\tau$ of that one-step law~\cite{bilmes1998,
grinsteadSnell1997}. The third estimator sets the emission parameters:
within each state with at least two observations and sample standard
deviation at least $10^{-12}\,\mathrm{yr}^{-1}$, $\mu_k$ and
$\sigma_k$ are the sample mean and standard deviation of the assigned
observations. Otherwise, the implementation uses the fallback rules
stated in the main text and written out in
Eq.~\eqref{eq:emission-fallback} of Online
Appendix~\ref{sec:supp-fallback}. In either case, $\sigma_k$ enters
Eq.~\eqref{eq:supp-emission-density} as the scale, so the emission
standard deviation is $\sigma_k \sqrt{\nu/(\nu - 2)}$. The law of
total variance in the next subsection quantifies the aggregate effect
of this scale convention (Eq.~\eqref{eq:supp-variance}).

The fourth estimator, the initial distribution, turns out to be
itself an empirical summary. Let $n_k$ count the observations
assigned to state $k$ over the $T$ observations, and let
$\hat\pi_k = n_k / T$ be the occupancy frequency. The number of
transitions out of state $j$ satisfies
$n_{j\cdot} = n_j - \mathbf{1}\{s_T = j\}$, and the number of
transitions into state $k$ satisfies
$n_{\cdot k} = n_k - \mathbf{1}\{s_1 = k\}$, because only the last
observation opens no transition and only the first closes none.
Substituting both identities, with italic $T$ the number of
observations, gives:
\begin{align}
  \big(\hat\pi\, \widehat{\mathbf{T}}\big)_k
   &\;=\; \sum_j \frac{n_j}{T} \cdot \frac{n_{jk}}{n_{j \cdot}}
     && \text{(definitions)} \notag \\
   &\;=\; \sum_j \frac{n_{j \cdot}}{T} \cdot \frac{n_{jk}}{n_{j \cdot}}
     \;+\; \sum_j \frac{\mathbf{1}\{s_T = j\}}{T}
       \cdot \frac{n_{jk}}{n_{j \cdot}}
     && \text{($n_j = n_{j\cdot} + \mathbf{1}\{s_T = j\}$)} \notag \\
   &\;=\; \frac{n_{\cdot k}}{T}
     \;+\; \frac{\widehat{\mathbf{T}}_{s_T, k}}{T}
     && \text{(sum transitions into $k$)} \notag \\
   &\;=\; \hat\pi_k
     \;+\; \frac{\widehat{\mathbf{T}}_{s_T, k}
       - \mathbf{1}\{s_1 = k\}}{T}.
     && \text{($n_{\cdot k} = n_k - \mathbf{1}\{s_1 = k\}$)}
  \label{eq:supp-fixed-point}
\end{align}
The occupancy frequencies therefore satisfy the stationarity equation
$\hat\pi \widehat{\mathbf{T}} = \hat\pi$ up to a fixed-point residual
bounded by $1/T$. Because the counted chain mixes within a few steps
(Online Appendix~\ref{sec:supp-gen-nojump}), a residual of this size
moves the solution by a comparable amount. For the fitted SPY model,
the solved eigenvector $\bar\pi$ deviates from the occupancy
frequencies by at most $3.7 \times 10^{-4}$, the same order as $1/T$,
against occupancy values
between $0.0058$ and $0.017$ (Table~\ref{tab:supp-generator-quantities}). The
initial distribution is consequently not an abstract spectral object:
$\bar\pi$ is, up to a correction of the order of $1/T$, the vector of
in-sample state frequencies.

\subsection{Stationary distribution and moments}
\label{sec:supp-gen-moments}

The stationary observation law of the no-jump generator follows
directly from Eq.~\eqref{eq:supp-joint-law}. Marginalizing the state
at a single time gives the mixture density and CDF:
\begin{equation}
  p(g) \;=\; \sum_{k=1}^{N} \bar\pi_k\, f_k(g),
  \qquad
  F_{\rm mix}(x) \;=\; \sum_{k=1}^{N} \bar\pi_k\,
  \Psi_\nu\!\left(\frac{x - \mu_k}{\sigma_k}\right),
  \label{eq:supp-mixture}
\end{equation}
where $\Psi_\nu$ is the CDF of the standard Student-$t$ with $\nu$
degrees of freedom. For SPY, every ingredient of
Eq.~\eqref{eq:supp-mixture} is a sample summary of the training
series: the component locations and scales are within-state sample
moments because no fallback was needed, and the stationary weights
differ from the occupancy frequencies only through the end effects of
Eq.~\eqref{eq:supp-fixed-point}. Weighting the empirical state means
by the occupancy frequencies regroups into the grand mean:
\begin{align}
  \sum_{k:n_k>0} \hat\pi_k\, \mu_k
   &\;=\; \sum_{k:n_k>0} \frac{n_k}{T} \cdot \frac{1}{n_k}
     \sum_{t \,:\, s_t = k} G_t
     && \text{(definitions of $\hat\pi_k$, $\mu_k$)} \notag \\
   &\;=\; \frac{1}{T} \sum_{t=1}^{T} G_t,
     && \text{(each $t$ counted once)}
  \label{eq:supp-mean}
\end{align}
and replacing $\hat\pi$ by $\bar\pi$ changes the value only through
the $1/T$ correction of Eq.~\eqref{eq:supp-fixed-point}: the model
mean is $0.0636~\mathrm{yr}^{-1}$ against the sample mean of
$0.0631~\mathrm{yr}^{-1}$ (Table~\ref{tab:supp-generator-quantities}).
For other assets, this identity needs a qualification. A state with a
single observation takes the full-series mean as its location
(Eq.~\eqref{eq:emission-fallback}), which need not preserve
Eq.~\eqref{eq:supp-mean}; the other fallback rule changes only the
scale (Table~\ref{tab:fallback_assets}).
The distributional match is equally direct.
Evaluated at the fitted SPY model, the largest gap between the model
CDF and the empirical CDF is
$\sup_x |F_{\rm mix}(x) - \widehat F_T(x)| = 0.0045$
(Table~\ref{tab:supp-generator-quantities}), where
$\widehat F_T$ is the in-sample empirical CDF. The path-level tests
compare a simulated path with the data, so the relevant yardstick is
the two-sample Kolmogorov--Smirnov (KS) critical value at significance
level $\alpha = 0.05$ for two samples of $2{,}766$ observations,
$1.358 \sqrt{2/T} = 0.0365$. The model-versus-data distance consumes
about one eighth of that threshold. Rejections of simulated paths are
therefore driven by path-level sampling variation rather than by a
misplaced marginal, which is consistent with the high in-sample
HMM-NJ pass rates in Table~\ref{tab:model_comparison}.

The second moment makes the role of the emission scale convention
explicit. Conditioning on the state and applying the law of total
variance gives:
\begin{align}
  \Var(G_t)
   &\;=\; \E\big[\Var(G_t \mid S_t)\big]
     \;+\; \Var\big(\E[G_t \mid S_t]\big)
     && \text{(law of total variance)} \notag \\
   &\;=\; \frac{\nu}{\nu - 2} \sum_k \bar\pi_k\, \sigma_k^2
     \;+\; \sum_k \bar\pi_k\, (\mu_k - \bar\mu)^2,
     && \text{(Student-$t$ variance)}
  \label{eq:supp-variance}
\end{align}
where $\bar\mu = \sum_k \bar\pi_k \mu_k$ is the stationary mean.
Define the within-state variance $V_{\rm w} = \sum_k \bar\pi_k
\sigma_k^2$ and the between-state variance $V_{\rm b} = \sum_k
\bar\pi_k (\mu_k - \bar\mu)^2$. The corresponding empirical
decomposition of the sample variance is $V_{\rm b} + V_{\rm w}$, up to
the $n_k - 1$ versus $n_k$ denominators of the within-state variances
and the $\bar\pi$ versus $\hat\pi$ weights, so the model variance
exceeds the sample variance by $2 V_{\rm w} / (\nu - 2)$. At the SPY fit
(Table~\ref{tab:supp-generator-quantities}),
$V_{\rm w} = 0.17~\mathrm{yr}^{-2}$, only $3.6\%$ of the total
variance of $4.60~\mathrm{yr}^{-2}$: the quantile partition
concentrates almost all dispersion between states rather than within
them. The inflation is correspondingly small, a model variance of
$4.72$ against $4.60~\mathrm{yr}^{-2}$ ($2.17$ against
$2.14~\mathrm{yr}^{-1}$ in standard deviation), and it is absent under
Gaussian emissions, whose model variance of $4.61~\mathrm{yr}^{-2}$
differs from the sample value only through those denominator and
weight corrections.

The fourth moment explains the emission comparison in
Table~\ref{tab:emission_hsmm}. Write $G_t = \mu_k +
\sigma_k X$ within state $k$, where $X$ is a standard Student-$t$
draw, and let $\delta_k = \mu_k - \bar\mu$. Expanding the fourth
central moment binomially and dropping the odd-power terms in $X$,
which vanish by the symmetry of $\psi_\nu$, gives:
\begin{align}
  \E\big[(G_t - \bar\mu)^4\big]
   &\;=\; \sum_k \bar\pi_k\,
     \E\big[(\delta_k + \sigma_k X)^4\big]
     && \text{(condition on the state)} \notag \\
   &\;=\; \sum_k \bar\pi_k \Big[
     \delta_k^4 + 4\, \delta_k^3 \sigma_k\, \E(X)
     + 6\, \delta_k^2 \sigma_k^2\, \E(X^2) \notag \\
   &\hspace{6em}
     + 4\, \delta_k \sigma_k^3\, \E(X^3) + \sigma_k^4\, \E(X^4)\Big]
     && \text{(binomial expansion)} \notag \\
   &\;=\; \sum_k \bar\pi_k \left[
     \delta_k^4
     + 6\, \delta_k^2\, \sigma_k^2\, \E(X^2)\right. \notag \\
   &\hspace{6em}\left.{}
     + \sigma_k^4\, \E(X^4)
     \right],
     && \text{($\E X = \E X^3 = 0$)}
  \label{eq:supp-fourth}
\end{align}
with the standard Student-$t$ moments, finite for $\nu > 4$:
\begin{equation}
  \E(X^2) \;=\; \frac{\nu}{\nu - 2},
  \qquad
  \E(X^4) \;=\; \frac{3 \nu^2}{(\nu - 2)(\nu - 4)}.
  \label{eq:supp-t-moments}
\end{equation}
At $\nu = 5$ these values are $5/3$ and $25$: the within-state
kurtosis is $9$ (excess kurtosis $6$), against $3$ (excess kurtosis
$0$) for Gaussian emissions. The excess kurtosis of the stationary mixture is
$\E[(G_t - \bar\mu)^4] / \Var(G_t)^2 - 3$; evaluating it at the SPY
fit gives $8.34$ with Student-$t$ emissions and $5.82$ with Gaussian
emissions, against an observed value of $7.72$
(Table~\ref{tab:supp-generator-quantities}). The emission choice enters
Eq.~\eqref{eq:supp-fourth} through both moments in
Eq.~\eqref{eq:supp-t-moments}; the larger change by far is
$\E(X^4)$ rising from $3$ to $25$, which acts on the $\sigma_k^4$
terms that the wide tail states dominate. The per-path simulated
HMM-WJ values in Table~\ref{tab:emission_hsmm} ($7.5$ and $5.5$) sit
below these no-jump population values for two reasons: jump episodes
lower kurtosis (Online Appendix~\ref{sec:supp-gen-tuning}), and the
sample kurtosis of a heavy-tailed series of $2{,}766$ observations is
biased downward, as the simulated in-sample HMM-NJ value of $8.1$
shows (Table~\ref{tab:model_comparison}). The population gap of $2.5$
between the emission choices and the simulated gap of $2.0$ have the
same direction and a similar size. The mixture also reproduces the asymmetry of the series without an
explicit skewness parameter: simulating $1{,}000$ in-sample-length
paths gives a mean sample skewness of $-0.76$ for HMM-NJ and $-0.73$
for HMM-WJ, against the observed $-0.75$
(Tables~\ref{tab:supp-generator-quantities}
and \ref{tab:descriptive_stats}). The fitted Laplace boundaries are
symmetric about the median, but the series is not, so the states on
the two sides of the median are filled unevenly: the within-state
means and scales and the state occupancies differ between the lower
and upper states, and the occupancy-weighted mixture inherits that
asymmetry.

\subsection{Dependence structure without jumps}
\label{sec:supp-gen-nojump}

The mixture and moment results describe the distribution at a single
date; the dependence over time is where the no-jump generator falls
short of the data. Within state $k$, the sojourn time is geometric:
the chain remains for $m$ steps with probability
$\mathbf{T}_{kk}^{\,m-1} (1 - \mathbf{T}_{kk})$ and therefore stays
for $1/(1 - \mathbf{T}_{kk})$ steps on average, between one and two
days for every fitted state
(Fig.~\ref{fig:model_internals}c). Short sojourns alone would not
erase the memory of return magnitude if the chain moved mostly to
neighboring states, but the counted transitions spread across the
whole range of states (Fig.~\ref{fig:model_internals}b), so the chain
soon loses track of the state it started in. The autocorrelation of
absolute returns makes this precise. Let $a_k =
\E\big[\,|G_t| \,\big|\, S_t = k\big]$ be the within-state mean
absolute return and $\bar a = \sum_k \bar\pi_k a_k$ its stationary
average. Because observations are conditionally independent given the
states, the lag-$\tau$ product moment reduces to a state-pair sum:
\begin{align}
  \E\big[\,|G_t|\, |G_{t+\tau}|\,\big]
   &\;=\; \sum_{j, k}
     \Pr(S_t = j,\, S_{t+\tau} = k)\; a_j\, a_k
     && \text{(conditional independence)} \notag \\
   &\;=\; \sum_{j, k} \bar\pi_j\,
     \big(\mathbf{T}^\tau\big)_{jk}\; a_j\, a_k,
     && \text{(Chapman--Kolmogorov)}
  \label{eq:supp-acf-nj}
\end{align}
so the autocorrelation function (ACF) of $|G_t|$ is:
\begin{equation}
  \mathrm{ACF}_{|G|}(\tau)
  \;=\; \frac{\sum_{j, k} \bar\pi_j\,
    \big(\mathbf{T}^\tau\big)_{jk}\, a_j\, a_k \;-\; \bar a^2}
    {\Var(|G_t|)}.
  \label{eq:supp-acf-nj-def}
\end{equation}
As $\tau$ grows, $\mathbf{T}^\tau$ converges
to its rank-one limit, every row equal to $\bar\pi$, at the geometric
rate $|\theta_2|^\tau$, where $\theta_2$ is the second-largest
eigenvalue of $\mathbf{T}$ in modulus; the ACF is therefore bounded
by a constant multiple of $|\theta_2|^\tau$. For the fitted SPY
matrix (Table~\ref{tab:supp-generator-quantities}), $|\theta_2| = 0.253$, and $|\theta_2|^\tau$ falls below
$0.01$ from lag $4$ onward. Evaluating Eq.~\eqref{eq:supp-acf-nj}
gives a model ACF of $0.26$ at lag $1$, close to the observed $0.30$
because the counted matrix pins the one-step state-pair law, but
$0.021$ at lag $3$ and effectively zero from lag $5$, while the
observed ACF remains near $0.26$ at lag $5$ and $0.14$ at lag $25$
(Table~\ref{tab:supp-generator-quantities};
Fig.~\ref{fig:empirical_motivation}d). For a first-order chain fitted
by transition counting, matching the one-step law fixes $\mathbf{T}$,
and the chain's memory is then set by that matrix's second
eigenvalue; for SPY, $|\theta_2| = 0.253$ forces mixing within about
a week. This is the quantitative content of the main-text statement
that the counted transition probabilities leave extreme-return states
too quickly to reproduce volatility clustering, and it is the failure
the jump mechanism is designed to address.

\subsection{Jump-episode arithmetic}
\label{sec:supp-gen-renewal}

The jump mechanism overlays the fast-mixing chain with rare, long
episodes, and its operating characteristics follow from renewal
arguments. On each step with no active episode, an episode with at
least one forced step begins with probability
$p_e = \epsilon (1 - e^{-\lambda})$, since a trigger whose Poisson
draw returns zero is discarded; conditional on starting, the episode
length is the zero-truncated Poisson draw with mean
$\E[K \mid K \ge 1] = \lambda / (1 - e^{-\lambda})$. At $\lambda =
90$, the correction $e^{-\lambda}$ is below $10^{-39}$, so $p_e =
\epsilon$ and $\E[K \mid K \ge 1] = \lambda$ to machine precision.
The long-run fraction of forced steps, a scalar denoted $\pi_J$,
follows from the renewal-reward theorem. A renewal cycle consists of
the free steps preceding an episode and the episode itself. The number
of steps up to and including the trigger is geometric with mean
$1/p_e$, but the triggering step is itself the first forced step, so
a cycle contains on average $1/p_e - 1 = (1 - p_e)/p_e$ free steps,
and the episode contributes $\E[K \mid K \ge 1]$ forced steps in
total:
\begin{align}
  \pi_J
   &\;=\; \frac{\E[K \mid K \ge 1]}
     {\dfrac{1 - p_e}{p_e} + \E[K \mid K \ge 1]}
     && \text{(forced steps over cycle length)} \notag \\
   &\;=\; \frac{p_e\, \E[K \mid K \ge 1]}
     {1 - p_e + p_e\, \E[K \mid K \ge 1]}
     && \text{(multiply through by $p_e$)} \notag \\
   &\;=\; \frac{\epsilon \lambda}
     {1 - \epsilon + \epsilon e^{-\lambda} + \epsilon \lambda}
     && \text{($p_e\, \E[K \mid K \ge 1] = \epsilon \lambda$)} \notag \\
   &\;\approx\; \frac{\epsilon \lambda}{1 + \epsilon \lambda}
     \;=\; 0.0089.
     && \text{(drop $\epsilon e^{-\lambda}$, then $\epsilon$)}
  \label{eq:supp-rho}
\end{align}
Fewer than one step in a hundred is forced at the fitted values; the
measured fraction across $400$ simulated in-sample-length paths is
$0.0088$ (Table~\ref{tab:supp-generator-quantities}). The product
$\epsilon \lambda$ sets this unconditional fraction, so the fraction
alone cannot separate the two parameters; the identification argument
of Online Appendix~\ref{sec:supp-gen-tuning} shows how conditioning on
paths with an episode separates them.

At the path level, the same arithmetic determines how episodes are
distributed across the ensemble. A simulated path contains no episode
only if every one of its $T - 1$ transition steps fails to trigger
one, so the jump-path probability at $T = 2{,}766$ is
$1 - (1 - p_e)^{T-1} = 24.2\%$; the simulated ensemble
gives $24.8\%$ ($99$ of $400$ paths;
Table~\ref{tab:supp-generator-quantities}), matching the $\sim\!24\%$
of paths with an episode that
Fig.~\ref{fig:statistical_validation} conditions on. We call these
paths jump paths. Dividing the $T - 1$ transition steps of a path by
the mean cycle length $(1 - p_e)/p_e + \lambda$ gives about $0.274$
episodes per path in expectation. The expected gap between episodes
is $1/\epsilon = 10^4$ trading days, about $40$ trading years, far
longer than the $11$-year path, so episodes are rare: roughly three quarters of the paths never trigger an
episode and follow the HMM-NJ law exactly, most of the remaining
quarter carry one episode, and about one jump path in eight carries
two or more. Within a jump path, the forced fraction is
$\lambda / T = 3.3\%$ for a single episode; the simulated mean is
$3.6\%$ (Table~\ref{tab:supp-generator-quantities}), slightly higher
because of the multi-episode paths.

The forced steps differ sharply from the free ones, which is what
drives the autocorrelation and kurtosis results of the next
subsection. Under the forced-step distribution $u$ of
Eq.~\eqref{eq:supp-u}, the conditional standard deviation is
$5.38~\mathrm{yr}^{-1}$ against the free-step $2.17~\mathrm{yr}^{-1}$,
a variance ratio of $6.1$, and the mean absolute return is
$\bar a_J = \sum_k u_k a_k = 4.96~\mathrm{yr}^{-1}$ against
$\bar a = 1.47~\mathrm{yr}^{-1}$
(Table~\ref{tab:supp-generator-quantities}). The same split describes
the state distribution of the augmented chain. On free steps the state
is distributed as $\bar\pi$, except during the few steps after an
episode ends, when the chain relaxes from the last forced tail state
back to $\bar\pi$ at the rate $|\theta_2|$. The long-run state
distribution is therefore approximately $(1 - \pi_J)\,\bar\pi +
\pi_J\, u$, and using $\bar\pi$ as the initial law misallocates
$O(\pi_J) \approx 1\%$ of probability mass. This quantifies the
main-text remark that $\bar\pi$ is exact for the ordinary transition
matrix but approximate for the jump-augmented process.

\subsection{Autocorrelation and kurtosis under jumps}
\label{sec:supp-gen-tuning}

The tuning objective of Eq.~\eqref{eq:grid_search} scores the
absolute-return ACF and the excess kurtosis, and both statistics
admit closed-form approximations that show what $(\epsilon, \lambda)$
control. The jump component of the ACF rests on one more renewal
quantity: the probability $h(\tau)$ that, given a step is forced, the
step $\tau$ later belongs to the same episode. A uniformly chosen
forced step falls in an episode of length $m$ with probability
proportional to $m \Pr(K = m \mid K \ge 1)$, because an episode of
length $m$ contains $m$ forced steps, and its position within that
episode is uniform. Combining the two gives:
\begin{align}
  h(\tau)
   &\;=\; \sum_{m \ge 1}
     \frac{m\, \Pr(K = m \mid K \ge 1)}{\E[K \mid K \ge 1]}
     \cdot \frac{(m - \tau)_+}{m}
     \notag \\
   &\;=\; \frac{\E\big[(K - \tau)_+ \,\big|\, K \ge 1\big]}
     {\E[K \mid K \ge 1]}
     && \text{(sum over $m$)} \notag \\
   &\;\approx\; \Big(1 - \frac{\tau}{\lambda}\Big)_{\!+},
     && \text{($K$ concentrated near $\lambda$)}
  \label{eq:supp-h}
\end{align}
where $(x)_+ = \max(x, 0)$ and the approximation holds because the
Poisson length concentrates within a few multiples of
$\sqrt{\lambda} \approx 9.5$ around $\lambda$. The same-episode
probability therefore decays almost linearly from one to zero over
the horizon $\lambda$. The ACF follows by conditioning on the
forced-step indicator $Z_t = \mathbf{1}\{\text{step $t$ is forced}\}$.
Given the indicator path, three kinds of pairs behave differently.
Two forced observations are conditionally independent, because every
forced step's state is an independent draw from $u$. A forced
observation and a later free one are linked only through the
episode's final state, and that link decays at the rate
$|\theta_2|^\tau$ of Online Appendix~\ref{sec:supp-gen-nojump}. Two
free observations are linked through the ordinary chain, which decays
at the same rate. Beyond the few lags over which $|\theta_2|^\tau$ is
non-negligible, only the conditional means covary. Write
$A_t = |G_t|$ for the absolute return. The conditional mean is
$\E[A_t \mid Z_t] = \bar a + Z_t\, (\bar a_J - \bar a)$,
and for rare episodes the joint forced probability is
$\Pr(Z_t = 1, Z_{t+\tau} = 1) \approx \pi_J\, h(\tau)$, so:
\begin{align}
  \Cov\big(A_t, A_{t+\tau}\big)
   &\;\approx\;
     \Cov\big(\E[A_t \mid Z_t],\;
              \E[A_{t+\tau} \mid Z_{t+\tau}]\big)
     && \text{(independence given $Z$)} \notag \\
   &\;=\; (\bar a_J - \bar a)^2\,
     \Cov\big(Z_t,\, Z_{t+\tau}\big)
     && \text{(conditional mean in $Z$)} \notag \\
   &\;\approx\; (\bar a_J - \bar a)^2\, \pi_J\, h(\tau),
     && \text{(episodes rare)}
  \label{eq:supp-acf-wj}
\end{align}
so the jump mechanism adds an ACF component of height $\pi_J (\bar a_J
- \bar a)^2 / \Var(A_t)$ that decays almost linearly to zero at lag
$\lambda$, precisely the slow component the counted chain cannot
produce. For a jump path carrying a single episode, the probability
that a step is forced is $\lambda / T$ rather than $\pi_J$, because
the path contains about $\lambda$ forced steps among its $T$ steps.
Replacing $\pi_J$ by $\lambda / T$ and using the absolute-return
variance of such a path, $\Var(A_t) = 3.01$, gives a predicted height
of $0.13$ at short lags, the plateau entry of
Table~\ref{tab:supp-generator-quantities}.
Table~\ref{tab:supp-acf-prediction} compares
this prediction with the mean simulated jump-path ACF: the two differ
by $0.013$ at lag $5$ and by less than $0.01$ from lag $10$ onward. At lags $1$ to $3$ the
simulated values exceed the episode component ($0.35$ against $0.13$
at lag $1$; Table~\ref{tab:supp-acf-prediction}) because the free-phase geometric component of
Section~\ref{sec:supp-gen-nojump} adds on top before it dies out.

% --- Table: predicted vs simulated jump-path ACF ---
\begin{table}[tbp]
\centering
\caption{\textbf{Episode-overlap prediction of the jump-path
    absolute-return autocorrelation function (ACF) against
    simulation.} The prediction evaluates
    Eq.~\eqref{eq:supp-acf-wj} at the fitted SPY values with the
    single-episode forced fraction $\lambda/T$ in place of $\pi_J$
    and the same-episode probability $h(\tau)$ of
    Eq.~\eqref{eq:supp-h}. The simulated column is the mean ACF of
    $|G_t|$ over the $99$ jump-active paths among $400$ simulated
    in-sample-length paths. The free-phase component of
    Eq.~\eqref{eq:supp-acf-nj} is excluded from the prediction; it
    has decayed below $0.01$ from lag $5$ onward.}
\label{tab:supp-acf-prediction}
\small
\input{sections/tables/tableS_acf_prediction}
\end{table}

Jump episodes lowered excess kurtosis while strengthening
absolute-return autocorrelation
(Table~\ref{tab:supp-generator-quantities}). Treating a path with one
episode as a mixture with forced-step weight $\lambda/T$, and using
$u$ in place of $\bar\pi$ for that component in
Eq.~\eqref{eq:supp-fourth}, gave population excess kurtosis of $7.20$
against $8.34$ without jumps. Forced steps increased variance by
about $17\%$ by drawing more often from the existing tail states.
The fourth central moment did not increase enough to offset the
increase in squared variance, so standardized kurtosis fell. To test
this tradeoff empirically, we split $4{,}000$ training-length HMM-WJ
paths into jump-active and no-jump groups and varied their mixture
weight from zero to one. A larger jump-active share lowered ACF-MAE
and KS and AD pass rates. Excess kurtosis moved from above the
observed $7.72$ at zero share to below it at a share of one, while
the Hill upper-tail index remained stable. The forced steps changed
how often the existing tail emissions were sampled, not their tail
exponent, which was set by $\nu$. At the realized share of about one
quarter, mean path excess kurtosis was $7.5$, slightly below the
observed value (Table~\ref{tab:model_comparison}). The objective in
Eq.~\eqref{eq:grid_search} balanced the gain in autocorrelation
against the change in excess kurtosis.

The same expressions show how the two jump parameters are identified,
and where they are not. Within a single-episode jump path, the ACF
component of Eq.~\eqref{eq:supp-acf-wj} has height proportional to
$\lambda / T$ and decays linearly over $\lambda$ lags, so over the
$25$ lags of the objective both its height and its slope depend on
$\lambda$; the kurtosis shift depends on the same forced fraction.
Every statistic entering Eq.~\eqref{eq:grid_search} for such a path
is therefore controlled by $\lambda$. The parameter $\epsilon$ moves
the objective only through the incidence of jump paths, which the
conditioning on jump paths removes, and through the small fraction of
multi-episode paths. The objective therefore varies little with
$\epsilon$, and the tuned value is not sharply determined by the
data: the tuned $\epsilon = 10^{-4}$ lies at the lower edge of the
explored range for SPY and for all three cross-asset tickers, and the
tuned optima vary with the Monte Carlo sample (Online
Appendix~\ref{sec:supp-cross-asset}).

% =============================================================================
% S5 Transition matrix internals
% =============================================================================

\section{Transition matrix internals and partition diagnostics}
\label{sec:supp-transitions}

Figure~\ref{fig:model_internals} shows the fitted components of the
SPY model at $N=100$ that enter the derivation of Online
Appendix~\ref{sec:supp-generator}. Panel~(a) shows that the fitted
Laplace CDF closely follows the empirical CDF; the $99$ boundaries
$Q_k$ are equiprobable under the fitted Laplace, and the empirical
state counts range from $16$ to $47$
(Table~\ref{tab:supp-generator-quantities}). Panel~(b) shows the
counted transition matrix $\mathbf{T}$: most of its entries are zero,
and the nonzero entries are scattered across the whole range of
destination states rather than concentrated near the diagonal, so
consecutive states are only weakly related, as the second eigenvalue
$|\theta_2| = 0.253$ quantifies (Online
Appendix~\ref{sec:supp-gen-nojump}). Panel~(c) shows that under the
ordinary transitions every state, including the ten tail states, has
an expected residence time of only $1$--$2$ steps. The chain also
leaves each tail set as a whole within one to two steps on average,
whereas a jump episode holds it in the tail sets for $\lambda = 90$
steps on average, the dashed line in panel~(c). Because every state
held observations, empty states and their uniform fallback rows
played no part in these short visits. The short visits, set against
the slow decay of the observed absolute-return autocorrelation, still
$0.14$ at lag $25$ (Table~\ref{tab:supp-generator-quantities}),
motivated the jump mechanism.

\begin{figure}[H]
  \centering
  \includegraphics[width=\textwidth]{figs/supplement/FigS05-Model-Internals.pdf}
  \caption{\textbf{Fitted hidden Markov model with jumps (HMM-WJ)
    internals for SPY with $N = 100$ states.} Panel~(a) overlays the
    fitted Laplace cumulative distribution function (CDF) on the empirical
    CDF; the $99$ vertical lines mark the quantile boundaries $Q_k$,
    equiprobable under the fitted Laplace, that define the hidden
    states. Panel~(b) displays the estimated transition matrix
    $\mathbf{T}$ in $\log_{10}$ color scale, with zero entries at the
    bottom of the scale; the nonzero entries are scattered across all
    destination states rather than concentrated near the diagonal.
    Panel~(c) plots the expected residence time
    $1/(1 - \mathbf{T}_{kk})$ under the ordinary transitions for each
    state on a $\log_{10}$ scale, with the
    bottom tail set $\mathcal{S}_{-}$ (states $1$--$5$, red) and the
    top tail set $\mathcal{S}_{+}$ (states $96$--$100$, teal)
    highlighted. The dashed horizontal line marks the fitted
    mean jump duration $\lambda = 90$ steps, far longer than the
    residence times of one to two steps.}
  \label{fig:model_internals}
\end{figure}

\section{Emission distribution and semi-Markov comparison}
\label{sec:supp-emission-hsmm}

To compare the emission and duration specifications, we evaluated
three variants on SPY (Table~\ref{tab:emission_hsmm}). The first is a
hidden semi-Markov model (HSMM) baseline adapted from Bulla and Bulla
\cite{bullaBulla2006}. It partitions the series into eight
Laplace-quantile states, the count with the highest in-sample KS pass
rate among $3$, $4$, $5$, $6$, and $8$ states. Each state's dwell
time is negative binomial, fitted by the method of moments when the
state has at least five visits with overdispersed lengths, and
geometric otherwise, and its emissions are Student-$t$ with $\nu = 5$.
The other two variants are HMM-WJ with Normal emissions and HMM-WJ
with Student-$t$ emissions, both with $N = 100$ states. In sample,
both HMM-WJ variants had higher KS and AD pass rates than the HSMM,
but the HMM-WJ variants use $100$ states and the HSMM uses $8$, so
this comparison does not assign the pass-rate difference solely to
the duration model. Out of sample, the HSMM matched or exceeded both
HMM-WJ variants on KS and AD pass rates. On temporal fit, the HSMM's
explicit dwell times did not lower the in-sample ACF-MAE below the
HMM-NJ value (Table~\ref{tab:sensitivity_n}), whereas both HMM-WJ
variants did. Holding the rest of HMM-WJ fixed, changing the emission
from Normal to Student-$t$ ($\nu=5$) increased simulated excess
kurtosis from $5.5$ to $7.5$, against an observed value of $7.7$
(Table~\ref{tab:descriptive_stats}). In sample, the Student-$t$
emission also increased the KS and AD pass rates; out of sample, the
two emission choices gave pass rates within one percentage point of
each other, and ACF-MAE and the distributional distances changed
little. The emission choice therefore explained the kurtosis gain in
this controlled HMM-WJ comparison, while it changed temporal fit
little.

\begin{table}[H]
\centering
\caption{\textbf{Three-way comparison of emission distributions and
    duration specifications for SPY.} Standard
    errors in parentheses; bold marks the best value per row, which for
    excess kurtosis is the value closest to the observed one. The HSMM
    is the hidden semi-Markov baseline adapted from Bulla and Bulla
    (eight states, negative-binomial dwell times with a geometric
    fallback, and Student-$t$ emissions) \cite{bullaBulla2006}; HMM-WJ:
    hidden Markov model with jumps
    ($N = 100$). All entries are computed from $1{,}000$ simulated
    paths; pass rates use significance level $\alpha = 0.05$. Kolmogorov--Smirnov
    (KS), Anderson--Darling (AD), and mean absolute error in the
    autocorrelation function (ACF-MAE) are defined in
    Online Appendix~\ref{sec:supp-validation}. Autocorrelation error
    compares the observed curve with the ensemble-average curve,
    using lags $1$--$252$ in training and $1$--$248$ in the holdout.}
\label{tab:emission_hsmm}
\small
\begin{tabular}{@{}lccc@{}}
\toprule
\textbf{Metric} & \textbf{HSMM} & \textbf{HMM-WJ} & \textbf{HMM-WJ} \\
                & \textbf{(Student-$t$)} & \textbf{(Gaussian)} & \textbf{(Student-$t$)} \\ \midrule
States              & $8$ & $100$ & $100$ \\[4pt]
\multicolumn{4}{l}{\textit{In-sample ($2{,}766$ trading days, $2014$--$2024$)}} \\[2pt]
KS pass rate (\%)   &  $82.0$ ($1.2$)         &  $97.0$ ($0.5$)         & $\mathbf{98.3}$ ($0.4$) \\
AD pass rate (\%)   &  $42.5$ ($1.6$)         &  $90.5$ ($0.9$)         & $\mathbf{94.8}$ ($0.7$) \\
Excess kurtosis     &   $4.8$ ($0.08$)        &   $5.5$ ($0.03$)        & $\mathbf{7.5}$ ($0.09$) \\
ACF-MAE             & $0.059$ ($<\!0.001$)    & $\mathbf{0.052}$ ($<\!0.001$) & $0.053$ ($<\!0.001$) \\
Wasserstein-$1$     & $0.176$ ($0.001$)       & $0.101$ ($0.002$)       & $\mathbf{0.097}$ ($0.001$) \\
Hellinger           & $0.113$ ($<\!0.001$)    & $0.076$ ($<\!0.001$)   & $\mathbf{0.074}$ ($<\!0.001$) \\[4pt]
\multicolumn{4}{l}{\textit{Out-of-sample ($249$ trading days, $2025$)}} \\[2pt]
KS pass rate (\%)   & $\mathbf{96.2}$ ($0.6$) &  $95.0$ ($0.7$)         &  $95.4$ ($0.7$) \\
AD pass rate (\%)   & $\mathbf{96.7}$ ($0.6$) &  $95.9$ ($0.6$)         &  $95.3$ ($0.7$) \\
Excess kurtosis     &   $4.1$ ($0.13$)        &   $4.9$ ($0.08$)        & $\mathbf{6.2}$ ($0.15$) \\
ACF-MAE             & $0.042$ ($<\!0.001$)    & $\mathbf{0.040}$ ($<\!0.001$) & $\mathbf{0.040}$ ($<\!0.001$) \\
Wasserstein-$1$     & $0.287$ ($0.002$)       & $\mathbf{0.275}$ ($0.005$) & $\mathbf{0.275}$ ($0.005$) \\
Hellinger           & $0.239$ ($0.001$)       & $0.212$ ($0.001$)       & $\mathbf{0.210}$ ($0.001$) \\
\bottomrule
\end{tabular}
\end{table}

% =============================================================================
% S7 State-resolution sensitivity
% =============================================================================

\section{State-resolution sensitivity}
\label{sec:supp-sensitivity}

To test sensitivity to state resolution, we refitted the hidden
Markov models with and without jumps (HMM-WJ and HMM-NJ) to the
$2{,}766$-day SPY training history using
$N\in\{30,60,90,100,150,200\}$ states
(Table~\ref{tab:sensitivity_n}). Each refit kept five states in each
tail set and repeated the $(\epsilon,\lambda)$ search. The no-jump model's
Kolmogorov--Smirnov (KS) and Anderson--Darling (AD) pass rates changed
by less than one percentage point, and its absolute-return
autocorrelation error (ACF-MAE) was unchanged at the reported
precision. For HMM-WJ, ACF-MAE fell from $0.057$ at $N=30$ to
$0.045$ at $N=200$, while the AD pass rate fell from $97.1\%$ to
$82.8\%$; intermediate values were not monotone. Increasing $N$
narrowed the fixed-size tail sets to more extreme quantile bands and
left fewer observations for each state's emission estimates, while
retuning also changed the jump parameters. The sweep therefore
measured their combined effect, without isolating the contribution of
each change. Fallback rules kept sparsely populated states defined
(Table~\ref{tab:fallback_assets}). The primary choice $N=100$ kept
the KS pass rate above $98\%$ and the AD pass rate within a few
percentage points of its best value while capturing part of the
autocorrelation improvement. We also used $N=100$ for other assets,
but this SPY-only sweep did not establish a suitable state count for
them or for shorter histories with fewer observations per state.

% --- Table: state-resolution sensitivity ---
\begin{table}[H]
\centering
\caption{\textbf{Sensitivity of in-sample fit to the number of
    states $N$.} All entries are computed from $1{,}000$ simulated
    in-sample-length paths for SPY at significance level
    $\alpha = 0.05$, with the tail-set size $N_{\rm tail}$ and the
    remaining hyperparameters of Table~\ref{tab:hyperparams} held
    fixed and the $(\epsilon, \lambda)$ grid search re-run at each
    $N$ for the hidden Markov model with jumps (HMM-WJ). Standard
    errors in parentheses: binomial for pass rates, bootstrap
    ($B = 500$) for the mean absolute error of the ensemble-average
    absolute-return autocorrelation function (ACF-MAE), using lags $1$--$252$. HMM-NJ denotes the hidden
    Markov model with no jumps. Values at $N = 100$ reproduce the
    corresponding entries of Table~\ref{tab:model_comparison}.}
\label{tab:sensitivity_n}
\small
\begin{tabular}{@{}lcccc@{}}
\toprule
\textbf{Model} & $\boldsymbol{N}$ & \textbf{KS pass (\%)} &
  \textbf{AD pass (\%)} & \textbf{ACF-MAE} \\ \midrule
HMM-NJ & $30$  & $99.3$ ($0.3$) & $97.8$ ($0.5$) & $0.059$ ($<\!0.001$) \\
HMM-NJ & $60$  & $99.5$ ($0.2$) & $98.5$ ($0.4$) & $0.059$ ($<\!0.001$) \\
HMM-NJ & $90$  & $99.3$ ($0.3$) & $98.6$ ($0.4$) & $0.059$ ($<\!0.001$) \\
HMM-NJ & $100$ & $99.3$ ($0.3$) & $98.2$ ($0.4$) & $0.059$ ($<\!0.001$) \\
HMM-NJ & $150$ & $99.1$ ($0.3$) & $98.7$ ($0.4$) & $0.059$ ($<\!0.001$) \\
HMM-NJ & $200$ & $99.2$ ($0.3$) & $98.3$ ($0.4$) & $0.059$ ($<\!0.001$) \\ \midrule
HMM-WJ & $30$  & $98.4$ ($0.4$) & $97.1$ ($0.5$) & $0.057$ ($<\!0.001$) \\
HMM-WJ & $60$  & $96.7$ ($0.6$) & $91.4$ ($0.9$) & $0.053$ ($<\!0.001$) \\
HMM-WJ & $90$  & $98.3$ ($0.4$) & $95.2$ ($0.7$) & $0.055$ ($<\!0.001$) \\
HMM-WJ & $100$ & $98.3$ ($0.4$) & $94.8$ ($0.7$) & $0.053$ ($<\!0.001$) \\
HMM-WJ & $150$ & $97.7$ ($0.5$) & $91.7$ ($0.9$) & $0.052$ ($<\!0.001$) \\
HMM-WJ & $200$ & $96.0$ ($0.6$) & $82.8$ ($1.2$) & $0.045$ ($0.001$) \\
\bottomrule
\end{tabular}
\end{table}


% =============================================================================
% S8 Cross-asset generalization (NVDA, JNJ, JPM)
% =============================================================================

\section{Cross-asset generalization}
\label{sec:supp-cross-asset}

To test whether the single-asset procedure generalized beyond SPY, we
repeated the fitting and evaluation on three stocks: NVDA
(semiconductors, high volatility), JNJ (health care, low
volatility), and JPM (financials, moderate volatility). Per-ticker
tune outputs at $(n_{\mathrm{paths}} = 200, \mathrm{seed} = 1234)$
were reproducible with these settings (Table~\ref{tab:cross_asset}).
The objective combines
ACF and kurtosis errors with fixed weights. Because kurtosis varies
across finite samples of heavy-tailed paths, the selected grid point
can change with the Monte Carlo sample. The reported
$(\epsilon^{*},\lambda^{*})$ values are reproducible with the stated
path count and seed, but they are not unique per-ticker optima. At the
selected parameter values, in-sample HMM-WJ KS pass rates were
$98.9\%$ for NVDA, $99.4\%$ for JNJ, and $99.1\%$ for JPM.
Out-of-sample KS pass rates on the $249$-day $2025$ hold-out were
$97.7\%$, $74.0\%$, and $79.2\%$, respectively. Thus the strong
in-sample fit transferred unevenly across assets, especially on the
short holdout.

% --- Table: Cross-asset generalization ---
\begin{table}[H]
\centering
\caption{\textbf{Cross-asset generalization of the hidden Markov model
    with jumps (HMM-WJ) across three
    single-stock tickers spanning sector and volatility regimes.}
    Per-ticker tune outputs $(\epsilon^{*}, \lambda^{*})$ at
    $n_{\mathrm{paths}} = 200$ and $\mathrm{seed} = 1234$, and the
    corresponding HMM-WJ Kolmogorov--Smirnov (KS) and Anderson--Darling
    (AD) pass rates on $1{,}000$ simulated paths at $\alpha = 0.05$
    against each ticker's in-sample (IS; $2014$--$2024$) and out-of-sample
    (OoS; $2025$) growth-rate history. Values in parentheses are binomial
    standard errors (SEs) for pass rates. The
    $(\epsilon^{*}, \lambda^{*})$ values are reported as the parameter
    values used for the IS / OoS metrics rather than as canonical
    per-ticker optima.}
\label{tab:cross_asset}
\small
\begin{tabular}{@{}lcccccc@{}}
\toprule
\textbf{Ticker} & \textbf{$\epsilon^{*}$} & \textbf{$\lambda^{*}$}
  & \textbf{IS KS (\%)} & \textbf{IS AD (\%)}
  & \textbf{OoS KS (\%)} & \textbf{OoS AD (\%)} \\ \midrule
NVDA (semiconductors)    & $10^{-4}$ & $70$ & $98.9$ ($0.3$) & $96.7$ ($0.6$) & $97.7$ ($0.5$) & $96.1$ ($0.6$) \\
JNJ (health care)        & $10^{-4}$ & $70$ & $99.4$ ($0.2$) & $96.3$ ($0.6$) & $74.0$ ($1.4$) & $61.3$ ($1.5$) \\
JPM (financials)         & $10^{-4}$ & $80$ & $99.1$ ($0.3$) & $95.7$ ($0.6$) & $79.2$ ($1.3$) & $77.7$ ($1.3$) \\
\bottomrule
\end{tabular}
\end{table}

\section{SIM composition derivations and properties verification}
\label{sec:supp-sim-derivations}

This appendix derives the variance scale of Eq.~\eqref{eq:scale} and
the residual-variance target of Eq.~\eqref{eq:r2-target}, and then the
mean, the large-$T$ OLS coefficients, and the composed variance that
result on each branch. Throughout, variances and covariances are
sample moments of the supplied paths; identities that drop the market
cross-covariance hold exactly when its sample value is zero and
approximately otherwise (Online Appendix~\ref{sec:supp-cov}). To
derive Eq.~\eqref{eq:scale}, first center the full-return
draw as in Eq.~\eqref{eq:center-generator}, then set
$\eps_i=s_i\tg_i^{\mathrm c}$ for $s_i\in[0,1]$ and require
that the composed series have the same variance as the generator
path:
\begin{equation}
  \Var(g_i) \;=\; \Var(\alpha_i + \beta_i\, g_m + \eps_i)
            \;=\; \Var(\beta_i\, g_m + \eps_i)
            \;=\; \sigma_{\gen,i}^2,
\end{equation}
where $\alpha_i$ drops out because it is constant. Expanding the sum
gives:
\begin{equation}
  \Var(\beta_i\, g_m + \eps_i) \;=\; \beta_i^2\, \Var(g_m)
  \;+\; 2 \beta_i\, \Cov(g_m,\, \eps_i)
  \;+\; \Var(\eps_i).
\end{equation}
Substituting $\eps_i = s_i\,\tg_i^{\mathrm c}$ gives
$\Var(\eps_i) = s_i^2\,\sigma_{\gen,i}^2$ and
$\Cov(g_m,\, \eps_i) = s_i\,\Cov(g_m,\, \tg_i^{\mathrm c})$. Under the
assumption $\Cov(\tg_i^{\mathrm c},\,g_m)\approx0$, the cross term vanishes and
the variance condition becomes
$\beta_i^2\, \sigma_m^2 + s_i^2\, \sigma_{\gen,i}^2 =
\sigma_{\gen,i}^2$. Solving for $s_i^2$ gives Eq.~\eqref{eq:scale}.

To derive Eq.~\eqref{eq:r2-target}, use the same zero-cross-covariance
assumption and write the $R^2$ of the composed series of
Eq.~\eqref{eq:sim-composition} regressed on the market:
\begin{equation}
  R^2 \;=\; \frac{\beta_i^2\, \sigma_m^2}{\beta_i^2\, \sigma_m^2 + \sigma_{\eps,i}^2}.
\end{equation}
Setting $R^2$ equal to the calibrated $R^2_{i,\real}$ and solving for
the residual variance gives the target $\sigma_{\eps,\mathrm{target}}^2$
of Eq.~\eqref{eq:r2-target}. This branch applies only to assets whose
calibrated $R^2$ exceeds $R^2_{\mathrm{preserve}} = 0.80$. In our
universe, only two index exchange-traded funds, QQQ and SPYG, cross
that threshold (Figs.~\ref{fig:r2_distribution}
and~\ref{fig:branch_map}), so the $R^2$-preserving branch composed two
assets and the variance-preserving branch composed the other $421$
(Table~\ref{tab:branch}).

% --- Figure: R^2 distribution across the universe ---
\begin{figure}[H]
  \centering
  \includegraphics[width=0.85\textwidth]{figs/supplement/FigS03-R2-Distribution.pdf}
  \caption{\textbf{Calibrated $R^2$ distribution across the $423$
    non-SPY tickers.} Histogram of ordinary least-squares (OLS) $R^2$
    against SPY over the
    $2014$--$2024$ window. The two tickers crossing the
    $R^2_{\mathrm{preserve}} = 0.80$ threshold (dashed line) are index
    exchange-traded funds (QQQ, SPYG); the highest-$R^2$ individual
    stock (BLK, $0.645$)
    falls $\sim 0.16$ below the threshold. The $R^2$-preserving
    branch therefore contains only QQQ and SPYG in this universe.}
  \label{fig:r2_distribution}
\end{figure}

% --- Figure: Branch map in (beta, R^2) space ---
\begin{figure}[H]
  \centering
  \includegraphics[width=0.85\textwidth]{figs/supplement/FigS06-Branch-Map.pdf}
  \caption{\textbf{Branch assignment in calibrated
    $(\beta, R^2_{\real})$ space.}
    The $R^2_{\mathrm{preserve}} = 0.80$ threshold (dashed line) separates
    tracker assets from single stocks. In this universe, only
    QQQ and SPYG exceed the threshold, and both are assigned to the
    $R^2$-preserving branch. The remaining $421$ tickers populate the
    variance-preserving branch; the clipped variance-preserving
    branch is unused because the market loading ratio
    $\rho_i = \beta_i^2 \sigma_m^2 / \sigma_{\gen,i}^2$ stays below
    $0.90$ for every ticker.}
  \label{fig:branch_map}
\end{figure}

% --- Table: Hybrid composer broken out by branch ---
\begin{table}[H]
  \centering
  \caption{\textbf{Hybrid composition method broken out by branch.}
    $N_{\mathrm{tickers}}$: number of tickers assigned to each branch by
    the $R^2_{\mathrm{preserve}} = 0.80$ threshold in a training run on
    the observed SPY path ($100$ replications per ticker); error, $W_1$,
    and $\kappa$ columns are medians across paths. The row labeled
    \texttt{hybrid} is the variance-preserving branch; its clipped case
    never occurred, because $\rho_i$ stayed below $1 - f = 0.90$ on
    every composed path. The
    two trackers in the $R^2$-preserving branch (QQQ and SPYG) had
    median absolute errors of $0.005$ for $\beta$ and $0.001$ for
    $R^2$, with a KS pass rate of $100\%$.}
  \label{tab:branch}
  \small
  \input{sections/tables/table2_by_branch}
\end{table}

Because the $R^2$-preserving branch composed only two assets, we also
checked it on a synthetic grid with supplied Gaussian residuals
(Table~\ref{tab:tracker_grid}). For each target loading and $R^2$, the
reference series combined the observed $2{,}766$-day SPY path, with
intercept zero, and an independent Gaussian residual of variance
$\sigma_{\eps,\mathrm{target}}^2$ from Eq.~\eqref{eq:r2-target}. A
second Gaussian residual draw was passed directly to
Algorithm~\ref{alg:hybrid}, together with the target variance
$\sigma_{\eps,\mathrm{target}}^2$ in place of the draw's sample
variance as $\sigma_{\gen}^2$, so the scale in
Eq.~\eqref{eq:r2-rescale} was exactly one. Only the composed residual
was centered over the full path; the reference residual kept its
random sample mean. The two residuals were also independent draws, so
the reference and composed series differed on any finite path, and
the loading and $R^2$ recovered by regressing the composed series on
the SPY path could differ from their targets. No full-return hidden
Markov model was fitted, and no new market path was simulated. This
grid therefore checked that every cell was assigned to the
$R^2$-preserving branch and that the branch recovers the target
loading and $R^2$ with Gaussian residuals, without evaluating
full-return reuse or temporal fit. Its Kolmogorov--Smirnov pass rates
are nominal only: both series contain the same SPY path, and the
composed residual is centered, so the two samples are not the
independent samples that the two-sample test assumes.

\begin{table}[H]
\centering
\caption{\textbf{Restricted Gaussian check of the $R^2$-preserving
    branch.} Both series used the observed SPY path and independent
    Gaussian residual draws; the composer used unit residual scale and
    centered its draw. No full-return model was fitted. Entries are
    medians over $100$ replications per cell with cross-replication
    standard deviations in parentheses. Loading and $R^2$ were recovered
    by ordinary least squares. The nominal Kolmogorov--Smirnov (KS) pass
    rate is the percentage of comparisons with $p>0.05$; its reference
    calibration was not established for this design. Every cell used
    the $R^2$-preserving branch.}
\label{tab:tracker_grid}
\small
\begin{tabular}{@{}ccccc@{}}
\toprule
$\boldsymbol{\beta_{\mathrm{true}}}$ &
$\boldsymbol{R^2_{\mathrm{true}}}$ &
$\boldsymbol{\hat\beta}$ & $\boldsymbol{\hat R^2}$ &
\textbf{Nominal KS pass (\%)} \\ \midrule
$0.8$ & $0.80$ & $0.8006$ ($0.0072$) & $0.7998$ ($0.0044$) & $100.0$ \\
$0.8$ & $0.85$ & $0.8006$ ($0.0059$) & $0.8497$ ($0.0041$) & $100.0$ \\
$0.8$ & $0.90$ & $0.8004$ ($0.0048$) & $0.8999$ ($0.0028$) & $100.0$ \\
$0.8$ & $0.95$ & $0.7995$ ($0.0033$) & $0.9500$ ($0.0013$) & $100.0$ \\
$0.8$ & $0.99$ & $0.7999$ ($0.0016$) & $0.9901$ ($0.0003$) & $100.0$ \\ \midrule
$1.0$ & $0.80$ & $1.0003$ ($0.0089$) & $0.8017$ ($0.0052$) & $100.0$ \\
$1.0$ & $0.85$ & $0.9994$ ($0.0084$) & $0.8497$ ($0.0039$) & $100.0$ \\
$1.0$ & $0.90$ & $0.9999$ ($0.0065$) & $0.8999$ ($0.0026$) & $100.0$ \\
$1.0$ & $0.95$ & $1.0005$ ($0.0036$) & $0.9502$ ($0.0013$) & $100.0$ \\
$1.0$ & $0.99$ & $1.0000$ ($0.0018$) & $0.9900$ ($0.0003$) & $100.0$ \\ \midrule
$1.2$ & $0.80$ & $1.2007$ ($0.0126$) & $0.8001$ ($0.0048$) & $100.0$ \\
$1.2$ & $0.85$ & $1.2011$ ($0.0105$) & $0.8503$ ($0.0043$) & $100.0$ \\
$1.2$ & $0.90$ & $1.1996$ ($0.0081$) & $0.8994$ ($0.0027$) & $100.0$ \\
$1.2$ & $0.95$ & $1.1995$ ($0.0052$) & $0.9498$ ($0.0014$) & $100.0$ \\
$1.2$ & $0.99$ & $1.1998$ ($0.0025$) & $0.9900$ ($0.0003$) & $100.0$ \\
\bottomrule
\end{tabular}
\end{table}

To distinguish intercept retention from mean preservation, let
$\overline{g}_m$ be the sample mean of the supplied market path.
Centering and scaling give $\overline{\eps}_i=0$ exactly, and the
implementation retains the fitted intercept $\alpha_i$ in every
branch. The resulting mean and its change from the calibrated linear
mean on the same market path are:
\begin{align}
  \overline{g}_i
    &= \alpha_i+\beta_i^{\eff}\overline{g}_m, \notag\\
  \overline{g}_i-(\alpha_i+\beta_i\overline{g}_m)
    &= (\beta_i^{\eff}-\beta_i)\overline{g}_m.
  \label{eq:supp-clipped-mean}
\end{align}
Retaining the intercept therefore preserves this linear mean when the
loading is unchanged, and a clipped path has a different mean whenever
$\overline{g}_m\ne0$. Clipping shrinks the magnitude of the loading,
so with a positive loading and a positive market mean the composed
mean falls, while other sign combinations raise it. The variance
budget and residual floor determine the clipped loading, with no
intercept adjustment to offset this mean change
(Table~\ref{tab:stress}). In the training comparison, the OLS
calibration also implies
$\alpha_i=\overline{g}_{i,\real}-\beta_i\overline{g}_{m,\real}$,
where the bars with subscript $\real$ denote observed training means.
The difference from the observed asset mean then separates into a
change in market mean and a change in loading:
\begin{equation}
  \overline{g}_i-\overline{g}_{i,\real}
  = \beta_i(\overline{g}_m-\overline{g}_{m,\real})
    +(\beta_i^{\eff}-\beta_i)\overline{g}_m.
  \label{eq:supp-training-mean-shift}
\end{equation}
Using the observed training SPY path with $\beta_i^{\eff}=\beta_i$
therefore matches the observed asset mean exactly. The stress
transformation of Table~\ref{tab:stress} preserved the market mean, so
only clipping could move the composed mean away from the observed
asset mean there. For a supplied market path with a different mean,
the observed asset mean need not be retained even without clipping.
The small intercept errors of the naive and hybrid methods therefore
reflect this imposed mean rather than independent recovery of
$\alpha_i$; centering the other methods' residuals gives them the
same small errors (Table~\ref{tab:centering_control}).

% --- Table: Clipping-branch stress test ---
\begin{table}[H]
\centering
\caption{\textbf{Clipping-branch stress test under market-volatility
    scaling.} We used
    $g_m^{(\gamma)}(t)=\overline{g}_m+
    \gamma[g_m(t)-\overline{g}_m]$, with
    $\gamma\in\{1,2,3\}$ and the calibrated marginals held fixed,
    and re-evaluated the full hybrid composition on the $423$-ticker
    universe. The transformation preserved the market mean and
    multiplied its standard deviation by $\gamma$. The fitted
    intercepts were retained, so clipping could change asset means
    according to Eq.~\eqref{eq:supp-clipped-mean}.}
\label{tab:stress}
\small
\begin{tabular}{@{}rr@{}}
\toprule
\textbf{$\gamma$} & \textbf{Clipped (\%)} \\ \midrule
$1$ & $0.0$  \\
$2$ & $76.4$ \\
$3$ & $95.2$ \\
\bottomrule
\end{tabular}
\end{table}

Retaining $\alpha_i$ also does not make the OLS intercept fitted to a
finite composed path equal $\alpha_i$ exactly. Because the composed
path is $g_i = \alpha_i + \beta_i^{\eff} g_m + \eps_i$, the OLS slope
and intercept on a finite path are:
\begin{align}
  \hat\beta_i
    &\;=\; \frac{\Cov(\alpha_i + \beta_i^{\eff} g_m + \eps_i,\, g_m)}{\Var(g_m)}
    && \text{(OLS slope)} \notag \\
    &\;=\; \beta_i^{\eff} + \frac{\Cov(\eps_i,\, g_m)}{\Var(g_m)},
    && \text{($\Cov(g_m, g_m) = \Var(g_m)$)} \notag \\
  \hat\alpha_i
    &\;=\; \overline{g}_i-\hat\beta_i\overline{g}_m
    && \text{(OLS intercept)} \notag \\
    &\;=\; \alpha_i-(\hat\beta_i-\beta_i^{\eff})\overline{g}_m.
    && \text{(Eq.~\eqref{eq:supp-clipped-mean})}
\end{align}
If $\eps_i$ is independent of $g_m$, the sample cross-covariance
vanishes as $T$ grows, so $\hat\beta_i \to \beta_i^{\eff}$ and
$\hat\alpha_i \to \alpha_i$. Here $\beta_i^{\eff}=\beta_i$ unless
Eq.~\eqref{eq:clip} clips the market loading. On the
variance-preserving branch, the composed variance equals
$\sigma_{\gen,i}^2$ when the sample cross-covariance is zero: without
clipping, $\beta_i^2 \sigma_m^2 + s_i^2 \sigma_{\gen,i}^2 =
\sigma_{\gen,i}^2$ by Eq.~\eqref{eq:scale}, and under clipping,
Eq.~\eqref{eq:clip} sets $\beta_i^{\eff}$ to give the same variance.
On the $R^2$-preserving branch, the composed variance is
$\beta_i^2 \sigma_m^2 / R^2_{i,\real}$, the value implied by the
calibrated $(\beta_i, R^2_{i,\real})$ rather than by the generator
marginal. On the variance-preserving branch, clipped or not, the
residual floor guarantees only that $\Var(\eps_i)\geq
f\,\sigma_{\gen,i}^2$; it does not guarantee that adding the market
term preserves the generator's tail distribution (Online
Appendix~\ref{sec:supp-residual-choice}).

The zero-sum constraint also fixes the residual contribution to the
terminal cumulative log return. With time step $\Delta t$, and because
the centered residuals sum to zero, summing the composed growth rates
gives the exact identity:
\begin{equation}
  \Delta t\sum_{t=1}^{T} g_i(t)
  = T\Delta t\,\alpha_i
    + \beta_i^{\eff}\Delta t\sum_{t=1}^{T}g_m(t),
  \qquad \sum_{t=1}^{T}\eps_i(t)=0.
  \label{eq:supp-zero-sum-terminal}
\end{equation}
For a fixed market path and effective loading, the residual contributes
no terminal uncertainty. Partial sums before $T$ can still vary. If
clipping changes the loading between draws, terminal returns can vary
through $\beta_i^{\eff}$, but the residual sum remains zero. Centering
uses the full horizon and is not equivalent to subtracting a fitted
population mean: if the same draw is extended to a longer horizon, its
first $T$ days receive a different centering constant and scale. These
paths are therefore constrained over a specified horizon and do not
represent unconstrained cumulative residual returns.

\subsection{Role of the centered full-return draw as the residual}
\label{sec:supp-residual-choice}

The hybrid composition places $\eps_i = s_i \tg_i^{\mathrm c}$, a
centered and rescaled draw of asset $i$'s full return, in the residual
slot of Eq.~\eqref{eq:sim-composition}. On its face this is a category
mismatch: the residual of a single-index decomposition is the
idiosyncratic component of the return, orthogonal to the market, while
$\tg_i^{\mathrm c}$ is distributed like the centered total return. Its
variance and tails include the part that, in observed data, comes
from co-movement with the market, although the draw itself is
simulated independently of the market path. This subsection states
what the substitution preserves exactly, what it preserves
approximately, and how its kurtosis changes as the market loading
grows.

The residual slot must supply three things. The first is a zero mean,
so that the intercept and market term alone set the conditional
mean. The second is the variance not already contributed by the
market term, $\sigma_{\gen,i}^2 - \beta_i^2 \sigma_m^2$ on the
unclipped variance-preserving branch. The third is the asset-specific
distributional shape: the heavy tails and volatility clustering that
motivated the single-asset generator. This subsection analyzes the
tails through kurtosis and does not analyze clustering. Centering
gives the residual path an exactly zero sample mean
(Eq.~\eqref{eq:center-generator}). On the unclipped variance-preserving
branch, computing $\sigma_m^2$ in Eq.~\eqref{eq:scale} from the same
market path used in the composition, as in training, gives the
residual exactly the remaining variance. The composed variance also
contains the sample cross-covariance term, which need not vanish on a
finite path (Online Appendix~\ref{sec:supp-cov}). Using the frozen
training market variance, as in the six-method holdout evaluation,
adds a further variance difference if the simulated market variance
differs from that training value. The full-return draw supplies the
residual's shape, but its validation as a standalone asset generator
does not establish that it reproduces the asset's regression-residual
distribution.

At $\rho_i = 0$, the loading is zero, Eq.~\eqref{eq:scale} gives
$s_i = 1$, and the composed path is the centered generator path
shifted to the mean $\alpha_i$. It therefore has the same centered
empirical distribution as the generator path, and it equals the
uncentered generator path only if $\alpha_i$ equals the draw's sample
mean. For $\rho_i > 0$, the independent market contribution changes
even the centered distribution. To describe that change on the
unclipped variance-preserving branch, we treat $s_i$ and $\beta_i$ as
constants rather than recomputing them from each draw, and we work
with population moments conditional on the fitted parameters, which
the sample moments of a long path approach. Write the deviation of the
composed return from its mean as $X_m + X_i$. The market contribution
is $X_m = \beta_i\, g_m^{\mathrm c}$, where $g_m^{\mathrm c}$ is the
market growth rate minus its mean, and the asset contribution is
$X_i = s_i\, \tg_i^{\mathrm c}$, the scaled asset draw. Their variances
are $v_m = \rho_i\, \sigma_{\gen,i}^2$ and
$v_i = (1 - \rho_i)\, \sigma_{\gen,i}^2$, and because the two are
independent, every cross term in the fourth moment factors:
\begin{align}
  \E\big[(X_m + X_i)^4\big]
   &\;=\; \E(X_m^4) + 4\,\E(X_m^3)\,\E(X_i)
     \notag \\
   &\qquad + 6\,\E(X_m^2)\,\E(X_i^2)
     + 4\,\E(X_m)\,\E(X_i^3)
     \notag \\
   &\qquad + \E(X_i^4)
     && \text{(independence)} \notag \\
   &\;=\; \E(X_m^4) + 6\, v_m\, v_i + \E(X_i^4).
     && \text{($\E X_m = \E X_i = 0$)}
  \label{eq:supp-comp-fourth}
\end{align}
Writing $\kappa_m$ and $\kappa_{\gen}$ for the excess kurtosis of the
supplied market series and of the generator draw, the component
fourth moments are $\E(X_m^4) = (\kappa_m + 3)\, v_m^2$ and
$\E(X_i^4) = (\kappa_{\gen} + 3)\, v_i^2$, where rescaling by $s_i$
leaves excess kurtosis unchanged. The excess kurtosis of the composed
return follows by dividing Eq.~\eqref{eq:supp-comp-fourth} by the
squared total variance:
\begin{align}
  \kappa(X_m + X_i)
   &\;=\; \frac{(\kappa_m + 3)\, v_m^2 + 6\, v_m\, v_i
     + (\kappa_{\gen} + 3)\, v_i^2}{(v_m + v_i)^2} \;-\; 3
     \notag \\
   &\;=\; \frac{\kappa_m\, v_m^2 + \kappa_{\gen}\, v_i^2
     + 3\, (v_m + v_i)^2}{(v_m + v_i)^2} \;-\; 3
     && \text{(collect terms)} \notag \\
   &\;=\; \kappa_m\, \rho_i^2 \;+\; \kappa_{\gen}\, (1 - \rho_i)^2.
     && \text{($v_m + v_i = \sigma_{\gen,i}^2$)}
  \label{eq:supp-kurt-comp}
\end{align}
The generator's excess kurtosis survives composition with weight
$(1 - \rho_i)^2$ and the market's enters with weight $\rho_i^2$; a
simulation check that combined pairs of independent paths from the
fitted SPY generator matched Eq.~\eqref{eq:supp-kurt-comp} to within
$0.03$ in mean excess kurtosis at $\rho_i \in \{0.1, 0.3, 0.5\}$
(Table~\ref{tab:supp-kurt-composition}). For a low-loading asset, most
of the generator's excess kurtosis survives: at $\rho_i = 0.10$, it
retains weight $0.81$. As $\rho_i$ grows, the composed shape moves
toward the market's, and the documented diagnostics show this
pattern. Median hybrid kurtosis falls from the lowest to the highest
$\beta$ quartile, as does naive kurtosis, and the hybrid KS pass rate
falls from $87.6\%$ to $65.9\%$ over the same range
(Table~\ref{tab:beta_bucket}). The Gaussian SIM, whose residual has no
excess kurtosis, stays far below both across the $\beta$ range
(Fig.~\ref{fig:tails}).

% --- Table: Per-beta-bucket breakdown ---
\begin{table}[H]
  \centering
  \caption{\textbf{Aggregate metrics by $\beta$ quartile, across composition methods.}
    Quartile cutoffs are computed from the calibrated $\beta$ distribution
    across the $423$-ticker universe. The hybrid KS pass rate declines
    from $87.6\%$ in Q1 to $65.9\%$ in Q4 as the market loading ratio
    $\rho_i$ grows; the naive KS pass rate falls
    from $24.6\%$ to $1.3\%$ over the same range.}
  \label{tab:beta_bucket}
  \small
  \input{sections/tables/table3_by_beta_bucket}
\end{table}

% --- Figure: Heavy-tail preservation ---
\begin{figure}[H]
  \centering
  \includegraphics[width=\textwidth]{figs/supplement/FigS02-Tail-Preservation.pdf}
  \caption{\textbf{Heavy-tail preservation: excess kurtosis and Hill
    tail index vs calibrated $\beta$.}
    Each point is a per-ticker median over $100$ replications.
    \emph{Left:} excess kurtosis, $y$-axis clipped to $[-2, 20]$ for
    readability. The dashed line marks the median real-data excess
    kurtosis across the universe ($\kappa_{\real} \approx 13$); the
    hybrid and naive composition methods (orange and blue) track each
    other, while the Gaussian single-index model (SIM) stays near
    $\kappa \approx 1$.
    \emph{Right:} Hill tail-index estimator on the upper $5\%$ of
    $|g|$. The hybrid and naive medians remain near $0.37$ throughout the
    $\beta$ range; Gaussian SIM sits near $0.22$. Because hybrid tracks naive
    across the $\beta$ range, the variance rescaling of the hybrid rule
    does not cause its shortfall from $\kappa_{\real}$; both carry the
    generator's tails, diluted by the market term as the market's share
    of the variance grows (Eq.~\eqref{eq:supp-kurt-comp}).}
  \label{fig:tails}
\end{figure}

The residual-fit methods offer a direct alternative: they fit a model
to each asset's OLS residual series (Table~\ref{tab:aggregate}). Those
residual series are orthogonal to the observed market by
construction, but a new simulated residual path need not have exactly
zero sample covariance with its market path. Fitting JumpHMM to
residuals also requires converting each residual series to a
synthetic price path by cumulative compounding and a second fit per
asset. The corrected full-return construction instead reuses the
generator's centered shape and supplies the residual variance budget
without that second fit. Its path mean equals
$\alpha_i+\beta_i^{\eff}\overline{g}_m$ exactly; matching the observed
training mean additionally requires the observed training market path
and an unclipped loading (Eq.~\eqref{eq:supp-training-mean-shift}).
Its variance and factor recovery depend on how $\sigma_m^2$ is
computed in Eq.~\eqref{eq:scale} and on the cross-covariance
assumption. At zero market loading it preserves the generator path's
centered empirical distribution, and for nonzero loading
Eq.~\eqref{eq:supp-kurt-comp} describes how the independent market
contribution changes the generator's excess kurtosis on the unclipped
variance-preserving branch. On the clipped and $R^2$-preserving
branches, the market's share of the composed variance differs from
$\rho_i$ (Algorithm~\ref{alg:hybrid}), and the same formula applies
with that share in place of $\rho_i$.

% --- Table: simulation check of the kurtosis composition rule ---
\begin{table}[H]
\centering
\caption{\textbf{Simulation check of the kurtosis composition rule.}
    For each variance share $\rho_i$, $200$ pairs of independent
    in-sample-length generator paths were standardized and combined
    as $\sqrt{\rho_i}\, x + \sqrt{1 - \rho_i}\, y$. The predicted
    column evaluates Eq.~\eqref{eq:supp-kurt-comp} with each pair's
    sample excess kurtosis; the simulated column is the sample excess
    kurtosis of the combined series. Both are averaged over the $200$
    pairs.}
\label{tab:supp-kurt-composition}
\small
\input{sections/tables/tableS_kurt_composition}
\end{table}


\section{Uncertainty in the multi-asset jump comparison}
\label{sec:supp-jump-composition}

The jump comparison used four seeds and $1{,}000$ market
replications per window, with all fits frozen (Table~\ref{tab:jump_ablation}). Within each setting,
naive and corrected composition shared the same asset and market
draws. The no-jump and market-only settings reused the no-jump asset
draws; the market-only and market-plus-asset settings reused the
jump-enabled market draws. Enabled and disabled models used separate
random-number streams, so matching replication identifiers did not
imply identical innovations. Primary summaries included every path;
jump-active subsets only described the mechanism.
For each metric, we averaged across tickers within each replication
and formed the difference between settings
(Table~\ref{tab:jump_ablation_uncertainty}). The Monte Carlo standard
error was the sample standard deviation of the $1{,}000$ differences
divided by $\sqrt{1{,}000}$, and intervals used the mean difference
plus or minus $1.96$ standard errors
(Table~\ref{tab:jump_ablation_uncertainty}). Keeping the cross-section
together retained dependence from the shared market draw. The
intervals did not treat assets as independent market histories or
include uncertainty from fitting or historical-window selection.
Kolmogorov--Smirnov (KS) non-rejection remained descriptive because
returns were temporally dependent.
The temporal contrast between training and holdout remained clear
in the replication-level uncertainty estimates
(Table~\ref{tab:jump_ablation_uncertainty}). The episode frequencies
matched the rare-event setting (Table~\ref{tab:jump_ablation_episodes}):
$247$ of
$1{,}000$ training-length SPY paths and $22$ of $1{,}000$ holdout-length
paths contained an episode. The corresponding shares across enabled
asset generators were $24.24\%$ and $2.46\%$. Market-only jumps
reduced average $25$-lag autocorrelation error for $99.5\%$ of
training assets but only $2.6\%$ of holdout assets. Enabling jumps in
both components reduced that error for $96.9\%$ and $0.5\%$,
respectively. These fractions described ticker-level mean changes,
not significance tests. The temporal contrast therefore extended
across most assets, within the limits of the short holdout and
transferred SPY calibration.

\begin{table}[H]
  \centering
  \caption{\textbf{Monte Carlo uncertainty in the effect of jump settings
    under corrected composition.} Differences are relative to jumps
    disabled, with $1{,}000$ market replications per window.
    $\Delta$ACF$_{25}$ is the change in mean absolute error of the
    absolute-return autocorrelation function through lag $25$;
    negative values indicate improved temporal fit.
    $\Delta$KS is the change in Kolmogorov--Smirnov non-rejection
    rate in percentage points. Brackets contain approximate $95\%$
    Monte Carlo (MC) intervals conditional on the fitted models and
    observed histories.}
  \label{tab:jump_ablation_uncertainty}
  \footnotesize
  \setlength{\tabcolsep}{4pt}
  \input{sections/tables/tableS_jump_ablation_uncertainty}
\end{table}

\begin{table}[H]
  \centering
  \caption{\textbf{Episode incidence and breadth of the temporal
    change in the multi-asset jump comparison.} SPY paths with an
    episode are counted over the $1{,}000$ jump-enabled market
    replications per window. Asset paths with an episode is the mean
    jump-active share across the enabled asset generators. The last
    two columns give the share of assets whose mean absolute-return
    autocorrelation error through lag $25$ fell under corrected
    composition when jumps were enabled in the market only, or in the
    market and the assets, relative to jumps disabled.}
  \label{tab:jump_ablation_episodes}
  \footnotesize
  \setlength{\tabcolsep}{4pt}
  \resizebox{\textwidth}{!}{\input{sections/tables/tableS_jump_ablation_episodes}}
\end{table}


\clearpage
\section{Residual-centering control}
\label{sec:supp-centering-control}

To separate the variance correction from the different residual mean
constraints, we paired each native training path with an equivalently
centered version (Table~\ref{tab:centering_control}). The experiment
reused the frozen full-return and residual models, ordinary least
squares (OLS) calibrations, observed training SPY, seed $1234$, and
$100$ draws per asset of length $2{,}766$. The Gaussian single-index
model (SIM), JumpHMM-on-residuals, stationary block bootstrap, and
generalized autoregressive conditional heteroskedasticity (GARCH)
model retained their native variance scales. With $g_i$ denoting a
native composed draw, the control removed only its residual mean:
\begin{equation}
  g_i^{\mathrm{ctrl}}(t)
  = g_i(t)-\frac{1}{T}\sum_{u=1}^{T}
    \left[g_i(u)-\alpha_i-\beta_i^{\eff}g_m(u)\right].
  \label{eq:supp-centering-control}
\end{equation}
The effective loading equaled the fitted loading except when the
hybrid rule changed it. Naive and hybrid paths already satisfied the
constraint and were reused without modification. No fits, random
draws, loadings, or scales changed between a native path and its
control. This paired translation left sample variance, recovered
loading, $R^2$, and excess kurtosis unchanged to numerical tolerance
for every evaluated path. The control assessed mean treatment; it did
not restore cumulative residual uncertainty or estimate parameter
uncertainty.

For a comparison on the same assets, the six-method summaries used
the $393$ assets whose GARCH fits passed the acceptance checks
(Table~\ref{tab:centering_control}). Refitting all $423$ residual
series with explicit convergence and admissibility checks reproduced
the accepted models' coefficients exactly. All optimizations converged;
the same $30$ fits were excluded because their ARCH and GARCH
coefficients summed to at least one. Acceptance records, coefficient
fingerprints, and exclusion reasons accompany the fitted-model cache.
The control therefore retained the same models and eligible asset set. There were $39{,}300$
paired paths per method on the common set. The accompanying results
also retain each method's full available population: $423$ assets
for the other five methods and $393$ for GARCH. All available-path
summaries reproduced the corresponding native values in
Table~\ref{tab:aggregate} at its reported precision for the metrics
included here. Pass rates pooled the same number of draws per asset;
Wasserstein distances and intercept errors were medians across paths.
Pass rates remained descriptive under serial dependence and the
fixed-sum constraint. Equivalent centering favored the residual-fit
comparators while preserving the improvement from naive to hybrid
composition, so the control narrowed the interpretation of the
cross-method comparison without changing the variance correction.

\begin{table}[H]
  \centering
  \caption{\textbf{Paired control for residual centering on common
    training assets.} Native and centered results use the same
    $393$ assets and $100$ paths per asset. Centered paths have zero
    residual mean over $2{,}766$ observations. Kolmogorov--Smirnov (KS)
    and Anderson--Darling (AD) columns give non-rejection percentages
    at $\alpha=0.05$; higher values indicate closer empirical fit.
    $W_1$ is median Wasserstein-1 distance, and $|\hat\alpha-\alpha|$
    is median absolute ordinary least squares (OLS) intercept error, both in
    $\mathrm{yr}^{-1}$; lower values indicate closer agreement.
    Naive and hybrid already center each residual path.
    JumpHMM denotes the hidden Markov generator, with jumps disabled;
    SIM denotes the single-index model, and GARCH denotes generalized
    autoregressive conditional heteroskedasticity.}
  \label{tab:centering_control}
  \footnotesize
  \setlength{\tabcolsep}{3pt}
  \input{sections/tables/tableS_centering_control}
\end{table}


\clearpage
\section{Pooled one-day tail thresholds}
\label{sec:supp-var-ensemble}

To estimate each generator's one-day tail without relying on the few
extreme observations in a single $249$-day path, we pooled returns
only after constructing separate paths
(Table~\ref{tab:var_convergence}). For every frozen training
asset and composition method, we generated $B=5{,}000$ paths of
length $249$ and calculated linearly interpolated empirical quantiles
from all $249B$ daily log returns. A market path was shared across
assets and methods within each replication. Each asset draw retained
its own centering and variance correction, with the observed training
market variance used in the hybrid rule. The hidden Markov models
started from their fitted stationary state distributions, jumps
remained disabled, and the generalized autoregressive conditional
heteroskedasticity (GARCH) simulator retained its $100$-step warmup.
We calibrated thresholds for all $423$ training assets, including
$393$ with accepted GARCH fits, before reading the holdout histories.
The resulting quantile estimated the return distribution of a
uniformly selected day within this finite-horizon construction.
It did not estimate a horizon-independent path law or restore
cumulative residual uncertainty.

To distinguish numerical stability from holdout performance, we fixed
pool sizes of $100$, $500$, $1{,}000$, $2{,}500$, and $5{,}000$ paths
before scoring the observed returns
(Table~\ref{tab:var_convergence}). An independent ensemble of
$1{,}000$ paths per asset and method checked the simulated probability
of falling below each threshold. Disjoint random streams separated
calibration and validation, while preserving the shared market within
each ensemble. Doubling the pool from $2{,}500$ to $5{,}000$ changed
the median asset's $99\%$ threshold by $0.12$--$0.30\%$, depending
on method. The corresponding $95$th percentiles of absolute relative
change were $0.37$--$1.03\%$. Individual thresholds were less stable
than aggregate rates: the largest $99\%$ change was $3.95\%$, for a
GARCH asset. Every method's mean observed exceedance rate changed by
less than $0.006$ percentage points at both coverage levels. Mean
exceedance rates on independent simulated paths were
$5.03$--$5.04\%$ and $1.01$--$1.02\%$, close to the respective
$5\%$ and $1\%$ targets. These checks assessed Monte Carlo stability
under the fitted generators; they did not establish coverage on real
returns. The reported cross-asset standard deviations were not
standard errors, and no coverage significance test treated shared
market histories or dependent days as independent observations.

\begin{table}[H]
  \centering
  \caption{\textbf{Threshold convergence and independent simulation
    checks.} Relative changes compare pooled thresholds from
    $2{,}500$ and $5{,}000$ separate $249$-day paths, divided by the
    absolute $5{,}000$-path threshold. Median and P95 denote the
    cross-asset median and $95$th percentile. Simulation rates are
    mean per-asset exceedance percentages on $1{,}000$ independent
    paths, using the $5{,}000$-path threshold. Summaries cover the
    $416$ holdout-eligible assets, or $386$ for generalized
    autoregressive conditional heteroskedasticity (GARCH).
    SIM denotes the single-index model; JumpHMM denotes the hidden
    Markov generator with jumps disabled.}
  \label{tab:var_convergence}
  \footnotesize
  \setlength{\tabcolsep}{3pt}
  \input{sections/tables/tableS_var_convergence}
\end{table}

To check whether differences in asset eligibility affected the
comparison, we restricted all six methods to the same $386$ assets
with accepted GARCH fits and complete holdout histories
(Table~\ref{tab:var_common}). The corrected method's $99\%$
exceedance rate was $1.08\%$, compared with $1.14\%$ for
JumpHMM-on-residuals and $0.85\%$ for naive composition. Gaussian
single-index composition remained farther above the nominal rate at
$1.66\%$. The common-asset comparison therefore retained the finding
that the corrected and residual-fit hidden Markov methods were close
to nominal on average. Neither the common-asset nor the full-asset
summary established equivalence, precise method rankings, or
calibration beyond the single observed market period.

\begin{table}[H]
  \centering
  \caption{\textbf{Pooled-threshold coverage on the common holdout
    assets.} All methods use the same $386$ assets. Each threshold
    pools $5{,}000$ separate $249$-day paths, as in
    Table~\ref{tab:var}. Rate denotes the mean per-asset exceedance
    percentage on observed $2025$ returns. Standard deviation (SD)
    describes cross-asset dispersion in percentage points.
    SIM denotes the single-index model; JumpHMM denotes the hidden
    Markov generator with jumps disabled; GARCH denotes generalized
    autoregressive conditional heteroskedasticity.}
  \label{tab:var_common}
  \footnotesize
  \setlength{\tabcolsep}{4pt}
  \input{sections/tables/tableS_var_common}
\end{table}


\section{Emission fallback inventory and sensitivity}
\label{sec:supp-fallback}

To identify where the fitted emissions differed from within-state
estimates, we reconstructed all state assignments, emission
parameters, and counted transition rows from the training inputs
(Table~\ref{tab:fallback_assets}). For one asset, let $n_k$ be the
number of observations in state $k$, let $\bar G_k$ and $s_k$ be
its sample mean and standard deviation when defined, and let
$\bar G$ and $s_G$ be the corresponding full-series quantities.
The implemented emission rule was:
\begin{equation}
 (\mu_k,\sigma_k)=
 \begin{cases}
   (\bar G,s_G), & n_k<2,\\
   (\bar G_k,s_G), & n_k\ge2,\ s_k<\tau,\\
   (\bar G_k,s_k), & n_k\ge2,\ s_k\ge\tau,
 \end{cases}
 \qquad \tau=10^{-12}\,\mathrm{yr}^{-1}.
 \label{eq:emission-fallback}
\end{equation}
The scale $\sigma_k$ multiplies a standard Student-$t$ draw; it is
not the emission standard deviation. State intervals included their
lower boundary and excluded their upper boundary. A transition row
with no observed outgoing transition was set to $1/N$ in every
column. Such a row was unconstrained by the observed transition
counts. Reconstruction matched every saved emission parameter and
transition probability exactly, including the corresponding fits to
regression residuals. Among the $424$ full-return models, $39$ had
at least one fallback emission, affecting $50$ states. Of these,
$39$ states contained $11$--$119$ identical observations each,
seven were empty, and four contained one observation. The seven
empty states also supplied the seven uniform transition rows.
Fallback states carried $0.399$--$4.300\%$ of stationary probability
per affected asset. Sparse-state probability was at most $0.0363\%$
per asset; most fallback mass belonged to populated states with
repeated returns. SPY required no fallback, and neither did any of
the $423$ residual-fit models.

\begin{table}[H]
 \centering
 \caption{\textbf{Fallback emissions in the full-return training fits.}
 Each listed asset had one populated constant-return state containing
 $n_{\rm c}$ observations. The columns $n_0$ and $n_1$ count empty
 and singleton states; the other $385$ full-return models required
 no fallback. Fallback mass is $100\sum_{k\in\mathcal F}\bar\pi_k$,
 where $\mathcal F$ contains all fallback states for that ticker.
 For populated constant states, the fitted location remained the
 state mean while the scale used the asset's full training standard
 deviation. Empty and singleton states used both full-series
 moments. $\Delta V$ is the percentage change in stationary mixture
 variance when only the populated constant-state scales are replaced
 by their empirical values, with all other parameters fixed.
 These are changes to the uncomposed generator, not variance errors
 from composition. All fits used $N=100$ and $\nu=5$.}
 \label{tab:fallback_assets}
 \footnotesize
 \setlength{\tabcolsep}{5pt}
 \input{sections/tables/tableS_fallback_assets}
\end{table}

To isolate the consequence of broadening populated constant states,
we compared the fitted generators with an empirical-scale alternative
on all $39$ affected non-market assets
(Table~\ref{tab:fallback_sensitivity}). The alternative replaced
$s_G$ by $s_k$ only where $n_k\ge2$ and $s_k<\tau$; locations,
partitions, transitions, initial probabilities, and sparse-state
fallbacks were unchanged. A zero empirical scale gives a point mass
at the state mean. The alternative was a diagnostic of how repeated
returns were represented, not a fitted or selected replacement model.
We replayed the original $100$ paths of $2{,}766$ observations per
asset, paired by state sequence and random draws, with jumps disabled
and the observed training SPY path shared by both treatments. Returns
from unmodified states were identical between treatments. We scored
the uncomposed generator and the naive and hybrid compositions;
each composition used its treatment's generator variance. All nine
original path metrics for the $7{,}800$ fitted naive and hybrid paths
matched the canonical cache exactly. Comparisons with the historical
returns used the original scoring functions, so pass fractions remain
descriptive summaries rather than calibrated significance claims for
dependent series containing ties.

The empirical-scale alternative improved marginal-fit scores on this
training subset (Table~\ref{tab:fallback_sensitivity}). For hybrid
composition, Kolmogorov--Smirnov (KS) pass rates increased from
$68.64\%$ to $78.92\%$, and Anderson--Darling (AD) pass rates
increased from $56.15\%$ to $70.23\%$. The paired changes were
$10.28$ and $14.08$ percentage points, with Monte Carlo standard
errors of $0.49$ and $0.60$ points. Mean Wasserstein distance fell
from $0.2782$ to $0.2474\,\mathrm{yr}^{-1}$. For the uncomposed
generators, the change in stationary mixture variance ranged from
$-0.651\%$ to $-6.520\%$ (Table~\ref{tab:fallback_assets}). The
fallback scale therefore affected both the generator's variance and
the composed marginal fit. All main comparisons retain the original
fallback rule. This paired training check did not test changes to
sparse-state estimates or validate the alternative on held-out
histories, and it did not establish a preferred production policy.

\begin{table}[H]
 \centering
 \caption{\textbf{Paired sensitivity to the scale of populated constant states.}
 Each treatment used $100$ paired training-length paths for each of
 the same $39$ affected assets. Entries are means over these fixed
 assets and replications. The empirical-scale treatment changes only
 populated near-constant states, retaining the sparse-state fallback.
 Kolmogorov--Smirnov (KS) and Anderson--Darling (AD) pass fractions
 use $p>0.05$; their changes are percentage points. $W_1$ is the
 Wasserstein-1 distance to observed growth rates. Variance / observed
 divides each path's variance by its asset's observed training
 variance. ACF-MAE$_{25}$ is the mean absolute difference between
 each path's absolute-return autocorrelation function and the observed
 function over lags $1$--$25$, calculated before averaging paths.
 Paired change subtracts the fitted value from the empirical-scale
 value. Standard errors (SE) use the $100$ replication-level mean
 paired changes over the fixed asset set and describe simulation
 uncertainty, not uncertainty across historical markets.}
 \label{tab:fallback_sensitivity}
 \footnotesize
 \setlength{\tabcolsep}{4pt}
 \input{sections/tables/tableS_fallback_sensitivity}
\end{table}
